\chapter{The three laws of thermodynamics}
\label{vol04:ch01}

\epigraph{The production of motive power is then due, in steam-engines,
not to an actual consumption of caloric, but to its transportation from
a warm body to a cold body, that is, to its re-establishment of
equilibrium.}{Sadi Carnot, \emph{Reflections on the Motive Power of
Fire} (1824), in the Mendoza translation \cite[p.~7]{hist:carnot1824}.}

\chapmeta{Half-life: foundational. Archetypes: balance, scaling. Exercise target: 38. Project track: household instrumentation. Hazard class: low (kitchen-scale temperatures and pressures). Reader path default: standard, with sections explicitly tagged. Named cases: none required in this chapter; perpetual-motion proposals are described by mechanism rather than by named example.}

\noindent
Thermodynamics is the engineering of energy accounting under
constraint. The first law fixes what is conserved; the second law fixes
which way the conserved quantity can move; the third law fixes the
zero. The chapter develops the three laws as bookkeeping rules,
identifies the state functions on which the bookkeeping depends, and
closes with the diagnostic for proposals that violate the rules. The
remaining chapters of the volume apply these rules to gases and
liquids (Ch 2), to engines and refrigerators (Ch 3), to entropy in its
microscopic form (Ch 4), and to heat and mass transfer (Ch 5-7).

\section{Energy as accounting}\pathtag{core}

A system is anything we draw a boundary around. The boundary may be
real (the metal wall of a calorimeter, the skin of a pipe) or
fictitious (a control volume in a flow, a thermodynamic envelope
around a building). Once the boundary is drawn, three transactions
cross it: matter, work, and heat. A system that exchanges all three
with its surroundings is an \engterm{open system}; a system that
exchanges work and heat but not matter is a \engterm{closed system}; a
system that exchanges nothing is an \engterm{isolated system}. The
distinction is set by the boundary, not by the contents.

A bank account is an accounting equation: deposits minus withdrawals
equals the change in balance. Thermodynamics is the same equation
applied to energy. The deposits are heat $Q$ added to the system and
work $W$ done on the system; the change in balance is the change
$\Delta E$ in the system's internal stored energy. The first law of
thermodynamics, in the form we will use throughout,
\[
\Delta E = Q + W,
\]
is the engineering accounting identity that every later chapter of
this volume builds on. The convention used here (work positive when
done on the system) is the one we will keep; some textbooks invert
the sign of $W$ \cite[ch.~2]{text:cengel-boles}. The chosen convention
is internally consistent; the reader who imports a result from another
text must check the sign on $W$ before using it.
Figure~\ref{fig:vol04ch01:first-law-signs} fixes the convention in a
single diagram.

% Figure: first-law sign conventions. A system boundary with Q in, W on the
% system, and mass flow arrows; the convention DeltaE = Q + W marked.
\begin{figure}[htbp]
\centering
\begin{tikzpicture}[>=latex, font=\small]
  % system box
  \draw[engblue, very thick, rounded corners] (0,0) rectangle (4,3);
  \node[engblue] at (2,2.4) {system};
  \node at (2,1.6) {internal energy $E$};
  \node at (2,1.0) {$\Delta E = Q + W$};
  % heat in (positive)
  \draw[->, engred, very thick] (-2,2.4) -- (0,2.4);
  \node[engred, anchor=south] at (-1,2.5) {$Q > 0$ (heat in)};
  % heat out (negative)
  \draw[->, engred, thick] (4,2.4) -- (6,2.4);
  \node[engred, anchor=south] at (5,2.5) {$Q < 0$ (heat out)};
  % work on the system (positive)
  \draw[->, englime!60!black, very thick] (-2,0.6) -- (0,0.6);
  \node[anchor=north, font=\scriptsize] at (-1,0.5) {$W > 0$ (work on system)};
  % work by the system (negative)
  \draw[->, englime!60!black, thick] (4,0.6) -- (6,0.6);
  \node[anchor=north, font=\scriptsize] at (5,0.5) {$W < 0$ (work by system)};
  % boundary label
  \node[anchor=north, font=\scriptsize, enggray] at (2,0) {boundary (real or fictitious)};
\end{tikzpicture}
\caption[First-law sign conventions]{Sign conventions used throughout this
volume. Heat $Q$ is positive when it crosses the boundary into the system;
work $W$ is positive when the surroundings do work on the system. With this
convention the closed-system balance reads $\Delta E = Q + W$. A gas that
expands and pushes a piston outward delivers work to the surroundings, so its
$W$ is negative. The reader who imports a formula from a text that puts work
done by the system as positive must flip the sign on every $W$ term.}
\label{fig:vol04ch01:first-law-signs}
\end{figure}


The accounting equation is exact. It is also unfalsified across two
centuries of experiment, despite many proposals to extract net work
from no input. The exactness has consequences:

\begin{itemize}[leftmargin=*]
  \item Energy is conserved across every closed boundary. A balance
    that fails to close points to an unaccounted flow, not to a
    failure of the law.
  \item Energy in transit has two engineering modes: heat (transfer
    driven by temperature difference) and work (transfer driven by
    force-times-displacement, voltage-times-charge, pressure-times-volume,
    chemical-potential-times-mass, or any equivalent product). The
    book keeps the two modes separate because the second law treats
    them asymmetrically.
  \item Internal energy is the cumulative state of the system; heat
    and work are not properties of the system, they are properties of
    a process the system undergoes. We will return to this distinction
    in section~\ref{sec:state-vs-path}.
\end{itemize}

The accounting identity has the same form for any control volume in
any domain. A battery, a steam turbine, a cell, a city, and a planet
all admit energy balances. The art of thermodynamic analysis is
choosing the boundary that makes the balance illuminating.

\begin{archetype}[Balance]
A balance equation states that the time rate of change of a conserved
quantity inside a control volume equals the net inflow minus the net
outflow, plus internal generation. Mass, momentum, energy, charge, and
species each obey balance equations of this form. The thermodynamic
balance is energy, the second-law balance is entropy (with a
generation term that cannot be negative), and Volume V will add
species balances for reacting systems. The reader who can write the
balance for any flow across any boundary owns the working core of the
engineering analysis of energy.
\end{archetype}

\section{First law: closed and open systems}\pathtag{core}

The first law for a \engterm{closed system} undergoing a process from
state 1 to state 2 is
\[
\Delta U = Q + W,
\]
where $U$ is the internal energy (a state function), $Q$ the heat
crossing the boundary into the system, and $W$ the work done on the
system. For a closed system at rest with no change in kinetic or
potential energy, $\Delta E = \Delta U$. When kinetic or potential
energy changes are non-negligible (an aircraft on take-off, a falling
body in a calorimeter), the full $\Delta E = \Delta U + \Delta KE +
\Delta PE$ is used.

\engterm{Boundary work} for a quasi-static process in a closed system
of fluid at pressure $p$ and volume $V$ is
\[
W = -\int_{V_1}^{V_2} p \,dV.
\]
The sign reflects our convention: when the system expands ($dV > 0$),
the surroundings do negative work on the system; the system delivers
work to the surroundings. The integral is path-dependent: the work
done by a gas in expanding from $(p_1, V_1)$ to $(p_2, V_2)$ depends
on the trajectory through the $p\textnormal{-}V$ plane, not only on
the endpoints. The four canonical quasi-static processes have
closed-form work integrals:

\begin{itemize}[leftmargin=*]
  \item Isobaric ($p$ constant): $W = -p (V_2 - V_1)$.
  \item Isochoric ($V$ constant): $W = 0$.
  \item Isothermal ideal gas: $W = -nRT \ln(V_2/V_1)$.
  \item Adiabatic ideal gas ($pV^\gamma$ constant): $W = (p_2 V_2 -
    p_1 V_1)/(\gamma - 1)$.
\end{itemize}

The four processes are the building blocks of every cycle in chapter~3
of this volume. Figure~\ref{fig:vol04ch01:pv-processes} draws all four
leaving a common start state, so the difference in the work (the area
under each curve) is visible at a glance.

% Figure: the four canonical quasi-static processes on a p-V diagram, all
% starting from the same state (V1=1, p1=4 arbitrary units).
\begin{figure}[htbp]
\centering
\begin{tikzpicture}
\begin{axis}[
  width=12cm, height=8cm,
  xlabel={volume $V$ (arb.)}, ylabel={pressure $p$ (arb.)},
  xmin=0.4, xmax=4.2, ymin=0, ymax=4.5,
  axis lines=left,
  legend pos=north east,
  legend cell align=left,
  tick label style={font=\scriptsize},
  label style={font=\small},
  legend style={font=\scriptsize},
  clip=false,
]
% common start state (V=1, p=4)
% isobaric: p constant = 4, V from 1 to 3
\addplot[engblue, very thick] coordinates {(1,4) (3,4)};
\addlegendentry{isobaric $p=\text{const}$}
% isochoric: V constant = 1, p from 4 to 1
\addplot[engred, very thick] coordinates {(1,4) (1,1)};
\addlegendentry{isochoric $V=\text{const}$}
% isothermal: pV = 4, from V=1 to V=4
\addplot[engorange, very thick, smooth, domain=1:4, samples=60] {4/x};
\addlegendentry{isothermal $pV=\text{const}$}
% adiabatic: pV^1.4 = 4, from V=1 to V=4 (steeper)
\addplot[englime!50!black, very thick, smooth, domain=1:4, samples=60] {4/(x^1.4)};
\addlegendentry{adiabatic $pV^{\gamma}=\text{const}$}
% common start marker
\filldraw[engblack] (axis cs:1,4) circle [radius=2pt];
\node[engblack, font=\scriptsize, anchor=south west] at (axis cs:1,4) {start};
\end{axis}
\end{tikzpicture}
\caption[The four canonical quasi-static processes]{The four canonical
quasi-static processes for an ideal gas, all leaving the same start state. The
boundary work is the area under each curve, $W = -\int p\,dV$, so the four
processes deliver different work for the same volume change. The isochoric path
does zero work (no volume change). The adiabatic curve falls faster than the
isothermal because the gas cools as it expands with no heat input, so its
pressure drops more steeply. These four segments are the building blocks of
every cycle in chapter 3 of this volume.}
\label{fig:vol04ch01:pv-processes}
\end{figure}


The script \texttt{code/ideal\_gas\_processes.py} evaluates $W$, $Q$,
and $\Delta S$ for each of the four processes from the same input state
and reproduces the worked numbers of section~\ref{sec:state-vs-path}.

The first law for an \engterm{open system} (also called a control
volume) includes mass flow. For a control volume at steady state with
inlet stream at enthalpy $h_1$ and outlet stream at $h_2$, mass flow
rate $\dot{m}$, heat rate $\dot{Q}$ added to the control volume, and
shaft work rate $\dot{W}_s$ done on the control volume,
\[
0 = \dot{Q} + \dot{W}_s + \dot{m}\left(h_1 - h_2 + \tfrac{1}{2}(V_1^2 -
V_2^2) + g(z_1 - z_2)\right).
\]
Figure~\ref{fig:vol04ch01:control-volume} lays out the terms.

% Figure: open-system (control volume) energy balance. One inlet stream,
% one outlet stream, heat in, shaft work, at steady state.
\begin{figure}[htbp]
\centering
\begin{tikzpicture}[>=latex, font=\small]
  % control volume
  \draw[engblue, very thick, dashed, rounded corners] (0,0) rectangle (5,3);
  \node[engblue, anchor=north] at (2.5,3) {control volume (steady state)};
  \node at (2.5,1.5) {$\dfrac{dE_{cv}}{dt} = 0$};
  % inlet stream
  \draw[->, engblack, very thick] (-2.2,2.2) -- (0,2.2);
  \node[anchor=south, font=\scriptsize] at (-1.1,2.3) {inlet $\dot m,\, h_1,\, V_1,\, z_1$};
  % outlet stream
  \draw[->, engblack, very thick] (5,2.2) -- (7.2,2.2);
  \node[anchor=south, font=\scriptsize] at (6.1,2.3) {outlet $\dot m,\, h_2,\, V_2,\, z_2$};
  % heat in
  \draw[->, engred, very thick] (1.0,-1.6) -- (1.0,0);
  \node[engred, anchor=north, font=\scriptsize] at (1.0,-1.6) {$\dot Q$};
  % shaft work
  \draw[->, englime!60!black, very thick] (4.0,-1.6) -- (4.0,0);
  \node[anchor=north, font=\scriptsize] at (4.0,-1.6) {$\dot W_s$};
  % balance equation
  \node[anchor=north, font=\scriptsize] at (2.5,-0.2)
    {$0 = \dot Q + \dot W_s + \dot m\!\left(h_1 - h_2 + \tfrac12(V_1^2-V_2^2) + g(z_1-z_2)\right)$};
\end{tikzpicture}
\caption[Open-system energy balance]{Steady-state energy balance on a control
volume with one inlet and one outlet. Mass enters carrying enthalpy $h_1$,
kinetic energy, and potential energy, and leaves carrying $h_2$ and its own
kinetic and potential energy; heat $\dot Q$ and shaft work $\dot W_s$ cross the
boundary. At steady state the stored energy inside the volume does not change,
so the inflows balance the outflows. The enthalpy $h = u + pv$ already contains
the flow work $pv$ needed to push mass across the boundary, which is why the
open-system balance is written in $h$ rather than $u$.}
\label{fig:vol04ch01:control-volume}
\end{figure}


The \engterm{enthalpy} $h = u + pv$ (per unit mass) groups the
internal energy with the flow work $pv$ required to push mass across
the control volume's boundary. Whenever mass crosses a boundary at
non-zero pressure, the relevant energy property is enthalpy, not
internal energy. The kinetic and potential terms are usually small for
process equipment (turbines, compressors, heat exchangers) and large
for nozzles, diffusers, and aircraft engines. The script
\texttt{code/first\_law\_balance.py} evaluates the residual of this
balance for any chosen control volume and solves for shaft power; it
is verified on the steam turbine of the Calculation exercises.

\engterm{Specific heats}. For a closed system at constant volume,
$dU = c_v(T) \,dT$ per unit mass, with $c_v$ the specific heat at
constant volume. For a closed system at constant pressure, $dH =
c_p(T) \,dT$, with $c_p$ the specific heat at constant pressure. The
two satisfy $c_p - c_v = R$ for an ideal gas (per mole) or $c_p -
c_v = R/M$ per unit mass with $M$ the molar mass. The ratio $\gamma =
c_p/c_v$ enters adiabatic processes and the speed of sound and is the
single most useful dimensionless property of a gas: $\gamma \approx
1.4$ for diatomic air at room temperature, $\gamma \approx 1.67$ for
monatomic argon, $\gamma \approx 1.3$ for steam at moderate
conditions.

\begin{principle}[The first law is an accounting identity]
The change in the energy of any thermodynamic system equals the net
energy crossing its boundary as heat, work, and (for open systems)
mass flow. The identity is exact, has units of joules (per unit time
for the rate form), and applies to any control volume in any domain.
Every later analysis of a thermal, electrical, chemical, or mechanical
process is the first law evaluated on a chosen boundary, with the
chosen work and heat terms identified. A balance that does not close
is a measurement or modelling error, not a violation of the law.
\end{principle}

\section{Second law: entropy and direction}\pathtag{core}

The first law says that energy is conserved. It does not say which way
energy is allowed to move. A flask of hot tea at 80~\degC{} placed in a
room at 20~\degC{} cools to 20~\degC{}; the room does not spontaneously
warm the tea to 80~\degC{}, even though both directions would conserve
energy. The asymmetry is the second law.

The two classical statements of the second law:

\begin{itemize}[leftmargin=*]
  \item \engterm{Clausius statement}: heat cannot, of itself, pass
    from a colder body to a hotter body. ``Of itself'' means without a
    compensating change elsewhere; refrigerators move heat from cold
    to hot, but only at the cost of input work that ends up as heat
    rejected at the hot side. The statement is the canonical
    expression of the directionality of heat flow
    \cite{hist:clausius1850}.
  \item \engterm{Kelvin-Planck statement}: no cyclic device can take
    heat from a single reservoir and convert it entirely to work. A
    cyclic heat engine must reject some heat to a colder reservoir;
    the ratio of work extracted to heat input is bounded by Carnot's
    theorem (chapter~3 of this volume).
\end{itemize}

The Kelvin-Planck statement is the schematic of
Figure~\ref{fig:vol04ch01:heat-engine}: a cyclic engine must reject
$Q_c$ to a cold reservoir, and the work it delivers is bounded by the
Carnot value.

% Figure: heat-reservoir / engine schematic. Hot reservoir Th, engine, cold
% reservoir Tc, with Qh in, W out, Qc rejected.
\begin{figure}[htbp]
\centering
\begin{tikzpicture}[>=latex, font=\small]
  % hot reservoir
  \draw[engred, very thick, fill=engred!12] (0,5) rectangle (5,6.2);
  \node[engred] at (2.5,5.6) {hot reservoir $T_h$};
  % cold reservoir
  \draw[engblue, very thick, fill=engblue!12] (0,0) rectangle (5,1.2);
  \node[engblue] at (2.5,0.6) {cold reservoir $T_c$};
  % engine
  \draw[engblack, very thick] (2.5,3.1) circle [radius=1.0];
  \node at (2.5,3.1) {engine};
  % Qh down into engine
  \draw[->, engred, very thick] (2.5,5.0) -- (2.5,4.1);
  \node[engred, anchor=west, font=\scriptsize] at (2.6,4.5) {$Q_h$};
  % Qc down to cold reservoir
  \draw[->, engblue, very thick] (2.5,2.1) -- (2.5,1.2);
  \node[engblue, anchor=west, font=\scriptsize] at (2.6,1.6) {$Q_c$};
  % W out
  \draw[->, englime!60!black, very thick] (3.5,3.1) -- (5.6,3.1);
  \node[anchor=south, font=\scriptsize] at (4.5,3.2) {$W = Q_h - Q_c$};
  % efficiency
  \node[anchor=west, font=\scriptsize] at (6.4,3.1) {$\eta = \dfrac{W}{Q_h} \le 1 - \dfrac{T_c}{T_h}$};
\end{tikzpicture}
\caption[Heat engine between two reservoirs]{A cyclic heat engine takes heat
$Q_h$ from a hot reservoir at $T_h$, delivers work $W$, and rejects heat
$Q_c = Q_h - W$ to a cold reservoir at $T_c$. The Kelvin-Planck statement of
the second law forbids $Q_c = 0$: no cyclic device turns heat from a single
reservoir entirely into work. The thermal efficiency $\eta = W/Q_h$ is bounded
above by the Carnot value $1 - T_c/T_h$, derived in chapter 3 of this volume.
A refrigerator or heat pump is the same diagram with every arrow reversed and
$W$ supplied rather than extracted.}
\label{fig:vol04ch01:heat-engine}
\end{figure}


The two statements are equivalent: a device that violated one could be
combined with a Carnot engine or refrigerator to produce a device that
violated the other. The proof is in Carnot's 1824 treatise in
embryonic form \cite{hist:carnot1824} and was made formal by Clausius
in 1850 \cite{hist:clausius1850}.

The bookkeeping form of the second law uses \engterm{entropy}, the
state function $S$ with
\[
dS = \frac{\delta Q_{\text{rev}}}{T},
\]
where $\delta Q_{\text{rev}}$ is the heat transferred to the system
along a reversible path and $T$ is the system's absolute temperature.
Entropy has units of joules per kelvin. For a closed system
undergoing any process between two states,
\[
\Delta S \geq \int \frac{\delta Q}{T},
\]
with equality for a reversible process and strict inequality for
irreversible processes. For an isolated system ($\delta Q = 0$), the
entropy can only increase or stay constant:
\[
\Delta S_{\text{isolated}} \geq 0.
\]
This is the form most often cited as the second law: the entropy of
the universe (taken as an isolated system) does not decrease.

The second law in rate form for an open system is the \engterm{entropy
balance}:
\[
\frac{dS_{cv}}{dt} = \sum \frac{\dot{Q}_j}{T_j} + \sum \dot{m}_i s_i -
\sum \dot{m}_e s_e + \dot{S}_{\text{gen}},
\]
where $T_j$ is the temperature at the boundary across which heat
$\dot{Q}_j$ enters, $s_i$ and $s_e$ are the inlet and exit specific
entropies, and $\dot{S}_{\text{gen}} \geq 0$ is the
\engterm{entropy-generation rate} inside the control volume. The
inequality $\dot{S}_{\text{gen}} \geq 0$ is the second law's
quantitative content for open systems: any process generates entropy
inside the control volume in an amount that depends on the process's
irreversibilities (heat transfer across a finite temperature
difference, friction, unrestrained expansion, mixing, chemical
reaction at finite affinity). Entropy generation is the engineering
diagnostic of loss; a process with $\dot{S}_{\text{gen}} = 0$ is
reversible and extracts the maximum possible work from its available
inputs. Real processes have $\dot{S}_{\text{gen}} > 0$ and extract
less.

The reader who finishes this chapter knows the second law as three
things at once: a direction-of-spontaneous-change rule, a bound on the
efficiency of heat engines, and an entropy balance with a
non-negative generation term. Chapter~4 of this volume provides the
microscopic interpretation that makes the entropy balance more than
bookkeeping.

\begin{mastery}
The bound $\dot{S}_{\text{gen}} \geq 0$ converts into a bound on lost
work through the \engterm{Gouy-Stodola theorem}. For a process drawing
on an environment at temperature $T_0$, the rate of lost work, the
shortfall between the actual work and the reversible best, is
\[
\dot{W}_{\text{lost}} = T_0 \, \dot{S}_{\text{gen}}.
\]
Every kelvin of entropy generated per second at a process with a
$300\,\si{\kelvin}$ environment destroys $300\,\text{W}$ of work
potential. This is the working content of second-law (exergy)
analysis: the engineer ranks the components of a plant by their
$T_0 \dot{S}_{\text{gen}}$ and attacks the largest term first. A heat
exchanger transferring $\dot{Q}$ across a temperature drop from $T_1$
to $T_2$ generates $\dot{S}_{\text{gen}} = \dot{Q}(1/T_2 - 1/T_1)$, so
the lost work scales with the temperature drop. Volume V develops the
exergy accounting; the present identity is the bridge from the entropy
balance to a euro figure.
\end{mastery}

\section{Third law}\pathtag{standard}

The third law of thermodynamics, in the form due to Walther Nernst
(1906) and made precise by Planck (1911), states that the entropy of a
pure crystalline substance at absolute zero is zero. The law has two
working consequences.

The first is the impossibility of reaching absolute zero in a finite
number of operations. Each cooling stage in a refrigeration cascade
extracts a finite amount of entropy; as $T \to 0$, the amount of
entropy that can be extracted per cycle approaches zero, and the
number of cycles required to reach $T = 0$ diverges. The lowest
temperatures achieved in laboratory experiments (nanokelvin in
laser-cooled atomic gases, picokelvin in nuclear-spin systems) are
finite for this reason.

The second is the use of absolute entropies. With the third law as
zero, the entropy of any substance at temperature $T$ above absolute
zero is obtained by integrating $c_p(T)/T \,dT$ from $0$ to $T$, plus
the latent heats divided by transition temperatures at any phase
transitions. Standard chemical and physical reference tables list
absolute molar entropies $S^\circ$ at the standard state ($25\,\degC$,
$1\,\text{atm}$) for thousands of substances. The absolute entropies
allow direct calculation of $\Delta S$ for chemical reactions from
tabulated values rather than from path integration.
Figure~\ref{fig:vol04ch01:entropy-temperature} shows the construction:
the curve climbs from the third-law zero by accumulating $c_p/T$, with
a vertical jump at each phase transition. A short table of standard
molar entropies is provided in \texttt{data/standard\_entropy.csv}.

% Figure: absolute entropy S(T) built from the third-law zero, integrating
% cp/T with jumps at phase transitions.
\begin{figure}[htbp]
\centering
\begin{tikzpicture}
\begin{axis}[
  width=11cm, height=7.5cm,
  xlabel={temperature $T$ (K)}, ylabel={absolute entropy $S$ (arb.)},
  xmin=0, xmax=400, ymin=0, ymax=6,
  axis lines=left,
  tick label style={font=\scriptsize},
  label style={font=\small},
  clip=false,
]
% solid branch: rises from (0,0), integral of cp/T, concave then linear-ish
\addplot[engblue, very thick, smooth] coordinates {
  (0,0) (20,0.25) (50,0.9) (100,1.7) (150,2.3) (200,2.75)
};
% melting jump at T=200: latent heat / T
\addplot[engred, very thick, dashed] coordinates {(200,2.75) (200,3.45)};
% liquid branch
\addplot[engblue, very thick, smooth] coordinates {
  (200,3.45) (250,3.9) (300,4.25) (350,4.5)
};
% vaporization jump at T=350
\addplot[engred, very thick, dashed] coordinates {(350,4.5) (350,5.6)};
% gas branch (short)
\addplot[engblue, very thick] coordinates {(350,5.6) (400,5.8)};
% annotations
\node[font=\scriptsize, anchor=west] at (axis cs:90,1.3) {solid};
\node[font=\scriptsize, anchor=west] at (axis cs:255,3.6) {liquid};
\node[font=\scriptsize, anchor=west] at (axis cs:360,5.6) {gas};
\node[engred, font=\scriptsize, anchor=west] at (axis cs:205,3.1) {fusion $L_f/T_m$};
\node[engred, font=\scriptsize, anchor=east] at (axis cs:345,5.1) {vaporization $L_v/T_b$};
\filldraw[engblack] (axis cs:0,0) circle [radius=2pt];
\node[font=\scriptsize, anchor=south west] at (axis cs:2,0) {$S(0)=0$};
\end{axis}
\end{tikzpicture}
\caption[Absolute entropy from the third-law zero]{The third law fixes the
entropy of a pure crystalline substance at absolute zero to be zero, which
turns entropy into an absolute quantity rather than a difference. The entropy
at any temperature is then $S(T) = \int_0^T (c_p/T)\,dT$ plus the latent heat
divided by the transition temperature at each phase change. The continuous
sloped segments are the solid, liquid, and gas branches; the vertical dashed
jumps are the entropy of fusion $L_f/T_m$ and of vaporization $L_v/T_b$.
Reference tables list the value reached at the standard state
($25\,\si{\degreeCelsius}$, $1\,\text{atm}$) as the standard molar entropy
$S^\circ$.}
\label{fig:vol04ch01:entropy-temperature}
\end{figure}


The third law is the least operationally useful of the three for
engineering practice at ordinary temperatures, but it is the law that
makes entropy a numerical quantity rather than a difference. In low-
temperature engineering (superconducting magnets, cryogenic liquids,
quantum computing hardware), the third law is the governing
constraint.

\begin{mastery}
The unreachability of absolute zero is quantified by adiabatic
demagnetisation, the workhorse of sub-kelvin cooling. A paramagnetic
salt is magnetised at a starting temperature while in thermal contact
with a bath, dumping its spin entropy; the bath is then removed and the
field lowered adiabatically (at constant entropy), which forces the
temperature down as the spins re-randomise. Each stage cools by a
factor set by the ratio of final to initial field, but the entropy
available to extract shrinks as $T \to 0$ because the spin system
freezes into its ground state, so the gain per stage collapses and no
finite sequence reaches $T = 0$. This is the operational content of the
Nernst statement: the isentropes of any substance converge as
$T \to 0$, and a process at constant entropy cannot bridge the gap to
the $T = 0$ isotherm in a finite number of steps. The same convergence
is why standard molar entropies are well defined: all substances share
the $S = 0$ floor.
\end{mastery}

\section{State functions vs path functions}\pathtag{core}
\label{sec:state-vs-path}

A \engterm{state function} is a property of the system that depends
only on the system's current state, not on how the system arrived
there. Internal energy $U$, enthalpy $H$, entropy $S$, Helmholtz free
energy $A$, Gibbs free energy $G$, temperature $T$, pressure $p$,
volume $V$, density $\rho$, and chemical composition are all state
functions. Their differentials are exact: $dU$, $dH$, $dS$ all
integrate around any closed cycle to zero.

A \engterm{path function} is a quantity associated with a process, not
with a state. Heat $Q$ and work $W$ are the canonical path functions:
the heat absorbed and the work done in going from state 1 to state 2
depend on the path through the $p$-$V$ plane (or any equivalent state
plane), not only on the endpoints. The differentials are inexact:
$\delta Q$ and $\delta W$ do not integrate to a state function around
a cycle. A heat engine extracts net work from a cyclic process
precisely because $W$ and $Q$ are path functions; $\oint dU = 0$
around any cycle, but $\oint \delta W \neq 0$ in general.

The distinction is the source of much of thermodynamics' analytical
power. The state-function property allows a calculation: to find
$\Delta H$ for a process, find any convenient reversible path between
the same endpoints and integrate. The path-function property forces a
real calculation: to find $W$ or $Q$, the actual process must be
modelled or measured.
Figure~\ref{fig:vol04ch01:state-vs-path} draws two paths between one
pair of states: the change in every state function is the same for
both, while the work (the area under each path) differs.

% Figure: two paths between the same two states on a p-V plane, illustrating
% that DeltaU is path-independent but W (area) differs.
\begin{figure}[htbp]
\centering
\begin{tikzpicture}
\begin{axis}[
  width=11cm, height=7.5cm,
  xlabel={volume $V$}, ylabel={pressure $p$},
  xmin=0.5, xmax=4.5, ymin=0, ymax=5,
  axis lines=left,
  legend pos=north east,
  legend cell align=left,
  tick label style={font=\scriptsize},
  label style={font=\small},
  legend style={font=\scriptsize},
  clip=false,
]
% State A = (1,4), State B = (4,1)
% Path I: isothermal pV=4 from A to B
\addplot[engblue, very thick, smooth, domain=1:4, samples=50] {4/x};
\addlegendentry{path I (isothermal)}
% Path II: down then across (isochoric then isobaric)
\addplot[engred, very thick] coordinates {(1,4) (1,1) (4,1)};
\addlegendentry{path II (isochoric + isobaric)}
% state markers
\filldraw[engblack] (axis cs:1,4) circle [radius=2.2pt];
\node[font=\scriptsize, anchor=south west] at (axis cs:1,4) {state $A$};
\filldraw[engblack] (axis cs:4,1) circle [radius=2.2pt];
\node[font=\scriptsize, anchor=south west] at (axis cs:4,1) {state $B$};
\end{axis}
\end{tikzpicture}
\caption[State functions versus path functions]{Two processes carry the system
from state $A$ to state $B$. The change in any state function, $\Delta U$,
$\Delta H$, $\Delta S$, is identical for both paths because it depends only on
the endpoints. The work $W = -\int p\,dV$ is the area under each curve, and the
two areas differ, so $W$ depends on the path taken. Heat $Q = \Delta U - W$
then also differs between the paths. Work and heat are path functions; the
state functions are not. This is why a cyclic engine can deliver net work even
though every state function returns to its starting value each cycle.}
\label{fig:vol04ch01:state-vs-path}
\end{figure}


\engterm{Common state-function identities for an ideal gas}. For an
ideal gas:
\begin{align*}
\Delta U &= n c_v \Delta T, \\
\Delta H &= n c_p \Delta T, \\
\Delta S &= n c_v \ln(T_2/T_1) + nR \ln(V_2/V_1), \\
\Delta S &= n c_p \ln(T_2/T_1) - nR \ln(p_2/p_1).
\end{align*}
The last two are the same equation in two parameter systems; they
follow from $dU = T\,dS - p\,dV$ (the combined first and second laws)
applied to an ideal gas with constant $c_v$. The reader who memorises
one form should be able to derive the other in under a minute.

\engterm{Cycle integrals}. Around a closed cycle the state functions
return to their initial values, so $\oint dU = 0$, $\oint dH = 0$,
$\oint dS = 0$. The first-law cycle integral $\oint \delta Q = -\oint
\delta W$ relates the net heat absorbed to the net work done; the
ratio (work to heat input) is the cycle's thermal efficiency, and the
Carnot bound (chapter~3) is the maximum achievable for given hot and
cold reservoir temperatures.

\section{Internal energy, enthalpy, free energies}\pathtag{standard}

Five state functions of energy underwrite the bulk of engineering
thermodynamic analysis. Each is defined by a combination of other
state functions, and each is the natural variable for a particular
process geometry.

\engterm{Internal energy} $U$. The total energy stored microscopically
in the system: kinetic energy of molecules, potential energy of
intermolecular forces, electronic energy, and the chemical-bond
energies. The natural variables of $U$ are $S$ and $V$:
\[
dU = T \,dS - p \,dV + \sum_i \mu_i \,dN_i,
\]
where $\mu_i$ is the chemical potential of species $i$ and $N_i$ the
number of moles. The chemical-potential term enters when the system's
composition can change (mass transfer, chemical reaction, phase
change).

\engterm{Enthalpy} $H = U + pV$. The natural variables of $H$ are $S$
and $p$:
\[
dH = T \,dS + V \,dp + \sum_i \mu_i \,dN_i.
\]
For a constant-pressure process, $\Delta H = Q$: enthalpy is the heat
exchanged by the system in any constant-pressure process. Combustion
processes, evaporation, and most chemical reactions in open laboratory
or industrial vessels are constant-pressure, so enthalpy is the
working quantity in chemical thermodynamics. Flow processes (where the
system is open and mass crosses a boundary) likewise use enthalpy
rather than internal energy because of the $pv$ flow-work term.

\engterm{Helmholtz free energy} $A = U - TS$. The natural variables of
$A$ are $T$ and $V$:
\[
dA = -S \,dT - p \,dV + \sum_i \mu_i \,dN_i.
\]
For a closed system at constant temperature and volume, $\Delta A$ is
the maximum useful work that can be extracted from a state change.
Constant-$T$ constant-$V$ processes are rare in engineering (the
typical closed-vessel reaction is at constant volume but variable
temperature), but $A$ is the natural function for statistical
mechanics: the canonical-ensemble partition function gives $A$
directly (chapter~4 of this volume).

\engterm{Gibbs free energy} $G = H - TS = U + pV - TS$. The natural
variables of $G$ are $T$ and $p$:
\[
dG = -S \,dT + V \,dp + \sum_i \mu_i \,dN_i.
\]
For a closed system at constant temperature and pressure, $\Delta G$
is the maximum useful work (other than $pV$ work against the
atmosphere) that can be extracted. Most chemical reactions of
engineering interest happen at approximately constant temperature and
pressure (open vessels at atmospheric, biological systems at body
temperature), so $G$ is the working free energy of chemistry and
biology. The equilibrium condition for a reaction at constant $T$, $p$
is $\Delta G = 0$; the criterion for spontaneity is $\Delta G < 0$;
the criterion for non-spontaneity is $\Delta G > 0$. Volume V develops
this machinery; Volume VI applies it to bioenergetics.

\engterm{Grand potential} $\Omega = U - TS - \sum_i \mu_i N_i$. The
natural variables are $T$, $V$, and $\mu_i$. The grand potential is
the working function for open-system statistical mechanics and for
problems where chemical potential rather than composition is the
externally controlled variable (membranes, electrochemistry,
semiconductor doping).

The five free energies are related by Legendre transforms; the
\engterm{Maxwell relations} that follow from the equality of cross
partial derivatives connect them and provide identities that are
indispensable in real-gas and phase-change calculations. Standard
examples:
\[
\left(\pdv{S}{V}\right)_T = \left(\pdv{p}{T}\right)_V, \quad
\left(\pdv{S}{p}\right)_T = -\left(\pdv{V}{T}\right)_p.
\]
The Maxwell relations let the engineer trade an experimentally
inaccessible derivative for one obtainable from $pVT$ data. Chapter~2
of this volume uses them to compute departures from ideal-gas
behaviour.

\section{Worked examples: gas compression, mixing, phase change}\pathtag{standard}

The state-function machinery becomes practical through repeated
application to canonical processes. Three examples cover the working
core.

\subsection{Isothermal and adiabatic compression of an ideal gas}

A piston compresses one mole of air (an ideal gas with $\gamma = 1.4$,
$R = 8.314\,\text{J}/(\text{mol}\cdot\text{K})$, $c_v = 5R/2$) from
$p_1 = 1\,\text{atm}$, $V_1 = 24\,\si{\liter}$ to half its initial
volume. The compression can be done isothermally (held at $T_1 =
293\,\si{\kelvin}$ by contact with a reservoir) or adiabatically
(quickly, with no heat transfer).

For the isothermal compression, $T$ is constant, so $\Delta U = 0$ and
$Q = -W$. The work done on the gas is $W = -nRT \ln(V_2/V_1) = -1
\times 8.314 \times 293 \times \ln(0.5) = +1690\,\text{J}$ (positive
because work is done on the gas; the convention is preserved). The
heat rejected to the reservoir is $-1690\,\text{J}$. The entropy
change of the gas is $\Delta S = nR \ln(V_2/V_1) = -5.76\,\text{J}/
\text{K}$, the reservoir's entropy change is $+5.76\,\text{J}/\text{K}$,
and the total entropy change is zero (a reversible isothermal
compression).

For the adiabatic compression, $Q = 0$ and $\Delta U = W$. The
end-state pressure and temperature follow from $pV^\gamma$ constant
and the ideal-gas law: $p_2 = p_1 (V_1/V_2)^\gamma = 2^{1.4} \times
1\,\text{atm} \approx 2.64\,\text{atm}$, $T_2 = T_1
(V_1/V_2)^{\gamma-1} = 2^{0.4} \times 293 \approx 387\,\si{\kelvin}$.
The work done on the gas is $W = n c_v \Delta T = 1 \times (5/2)
\times 8.314 \times (387 - 293) \approx 1960\,\text{J}$. The entropy
change of the gas is zero (reversible adiabatic), and the surroundings
exchange no heat, so the total entropy change is again zero.

The adiabatic compression takes more work than the isothermal because
the gas heats up; the increased pressure during compression demands
more work to push the piston the same distance. The same heating is
what makes a diesel engine's compression stroke ignite the fuel
charge: $T$ rises through the combustion temperature without an
electrical spark.

\subsection{Mixing of two ideal gases}

Two ideal gases at the same temperature and pressure, separated by a
partition, mix when the partition is removed. The internal energy of
the mixture equals the sum of the internal energies of the components
(an ideal-gas property), so $\Delta U = 0$ and $Q = 0$. No work is done
(the gases expand against each other, not against a piston), so $W =
0$ and the first-law balance is trivially $0 = 0 + 0$.

The entropy change is not zero. Each species expands into the full
volume; its partial pressure drops from the total pressure $p$ to its
partial pressure $x_i p$, where $x_i$ is its mole fraction. The
entropy change per mole of species $i$ is $-R \ln x_i$. For a mixture
of $n_A$ moles of $A$ and $n_B$ moles of $B$,
\[
\Delta S_{\text{mix}} = -R(n_A \ln x_A + n_B \ln x_B) > 0.
\]
The mixing entropy is positive because each $x_i < 1$ and the
logarithms are negative. The process is irreversible: separating the
two gases back into their original compartments requires work
(equivalent, in the ideal-gas limit, to the work of an isothermal
compression back to the original partial pressures). The script
\texttt{code/entropy\_of\_mixing.py} evaluates the formula and runs a
particle simulation whose bin-population entropy climbs to $R \ln 2$
per mole as an equimolar mixture equilibrates.

The mixing entropy is the engineering content of the second law for
separations. Distillation, gas absorption, membrane separation, and
crystallisation all undo a mixing-entropy increase, which costs work
in proportion to the entropy of the original mixing. Chapter~7 of
this volume develops the separation-process side; the present
calculation establishes the lower bound.

\subsection{Vapourisation of water}

A litre of water at $100\,\degC$ and $1\,\text{atm}$ is vapourised at
constant pressure. The latent heat of vapourisation of water at
$100\,\degC$ is $h_{fg} = 2257\,\text{kJ}/\text{kg}$ \cite[Table A-4]{
text:cengel-boles}. The water mass is $1\,\text{kg}$ (density
$1000\,\text{kg}/\text{m}^3$), so $Q = m h_{fg} = 2257\,\text{kJ}$.

The work done by the vapour against the atmosphere is $W = -p \Delta V
= -p(V_g - V_f)$. At $100\,\degC$ and $1\,\text{atm}$, the specific
volume of saturated vapour is $v_g \approx 1.673\,\text{m}^3/
\text{kg}$ and of saturated liquid $v_f \approx 1.04 \times
10^{-3}\,\text{m}^3/\text{kg}$, so $\Delta V \approx 1.67\,\text{m}^3
/\text{kg}$. The work is $W = -10^5 \,\text{Pa} \times 1.67\,
\text{m}^3 = -167\,\text{kJ}$ (negative: the system does work on the
atmosphere). The change in internal energy is $\Delta U = Q + W = 2257
- 167 = 2090\,\text{kJ}$.

The entropy change is $\Delta S = h_{fg}/T = 2257\,\text{kJ}/373\,
\si{\kelvin} = 6.05\,\text{kJ}/\text{K}$. The corresponding entropy
change of the surroundings (which gave up $2257\,\text{kJ}$ of heat at
the same temperature, idealised as reversible) is $-6.05\,\text{kJ}/
\text{K}$; the total entropy change is zero. A real kettle, of course,
delivers heat at a much higher temperature than the water's
$100\,\degC$; the actual surroundings entropy change is smaller in
magnitude, and the total entropy change is positive (irreversibility
due to finite-temperature-difference heat transfer). The saturation
properties used here are tabulated in
\texttt{data/water\_saturation.csv}.

\subsection{The Carnot cycle as an upper bound}

A cycle made of two isothermal and two reversible adiabatic segments,
the Carnot cycle, sets the efficiency ceiling for any heat engine
working between two reservoir temperatures. Take one mole of an ideal
gas with $\gamma = 1.4$ cycling between $T_h = 600\,\si{\kelvin}$ and
$T_c = 300\,\si{\kelvin}$, with an isothermal expansion that doubles
the volume at $T_h$. The heat absorbed in the high-temperature
expansion is $Q_h = nRT_h \ln 2 = 1 \times 8.314 \times 600 \times
0.693 = 3458\,\text{J}$. By the geometry of the cycle the entropy
delivered at the top equals the entropy removed at the bottom, so the
heat rejected is $Q_c = nRT_c \ln 2 = 1729\,\text{J}$, exactly half of
$Q_h$. The net work delivered is $Q_h - Q_c = 1729\,\text{J}$, and the
efficiency is
\[
\eta = \frac{W}{Q_h} = 1 - \frac{Q_c}{Q_h} = 1 - \frac{T_c}{T_h} =
1 - \frac{300}{600} = 0.50.
\]
The two heat exchanges happen at the single temperatures $T_h$ and
$T_c$, which is why the entropy transferred in and out are equal and
the efficiency reaches the bound. Any real engine between these two
reservoirs does worse, because real heat exchange happens across a
finite temperature difference and generates entropy.
Figures~\ref{fig:vol04ch01:carnot-pv} and~\ref{fig:vol04ch01:carnot-ts}
draw the cycle on $p$-$V$ and $T$-$S$ axes; on the second it is a
rectangle whose area is the net work. The script
\texttt{code/carnot\_cycle.py} computes the state points and confirms
that the cycle efficiency equals $1 - T_c/T_h$ to machine precision.

% Figure: the Carnot cycle on a p-V diagram. Two isotherms and two adiabats.
% State 1 (top): isothermal expansion to 2; adiabatic expansion to 3;
% isothermal compression to 4; adiabatic compression back to 1.
\begin{figure}[htbp]
\centering
\begin{tikzpicture}
\begin{axis}[
  width=11cm, height=8cm,
  xlabel={volume $V$ (arb.)}, ylabel={pressure $p$ (arb.)},
  xmin=0.8, xmax=6.5, ymin=0, ymax=6.2,
  axis lines=left,
  tick label style={font=\scriptsize},
  label style={font=\small},
  clip=false,
]
% Hot isotherm pV = 6 (T_h), cold isotherm pV = 2.4 (T_c).
% Adiabats pV^1.4 = const.
% State 1: V=1, p=6  (on hot isotherm, pV=6)
% State 2: V=2, p=3  (on hot isotherm, pV=6)
% State 3: on cold isotherm pV=2.4 and adiabat from 2: p2 V2^1.4 = 3*2^1.4=7.92
%   so p = 7.92/V^1.4 ; cold isotherm p=2.4/V => 2.4/V = 7.92/V^1.4
%   => V^0.4 = 7.92/2.4 = 3.3 => V = 3.3^2.5 = 19.8 too large; rescale cold isotherm.
% Use cold isotherm pV = 4.2: state 3 where 4.2/V = 7.92/V^1.4 => V^0.4=1.886 => V=4.0,p=1.05
% State 4 on cold isotherm and adiabat back to 1: p1 V1^1.4 = 6 ; p=6/V^1.4
%   cold isotherm p=4.2/V => 4.2/V = 6/V^1.4 => V^0.4 = 1.4286 => V=2.0, p=2.1
% Isothermal expansion 1->2 (hot, pV=6)
\addplot[engred, very thick, smooth, domain=1:2, samples=40] {6/x};
% Adiabatic expansion 2->3 (pV^1.4 = 7.92)
\addplot[englime!50!black, very thick, smooth, domain=2:4, samples=40] {7.92/(x^1.4)};
% Isothermal compression 3->4 (cold, pV=4.2)
\addplot[engblue, very thick, smooth, domain=2:4, samples=40] {4.2/x};
% Adiabatic compression 4->1 (pV^1.4 = 6)
\addplot[engorange, very thick, smooth, domain=1:2, samples=40] {6/(x^1.4)};
% state markers
\filldraw[engblack] (axis cs:1,6) circle [radius=2pt];
\node[font=\scriptsize, anchor=south west] at (axis cs:1,6) {1};
\filldraw[engblack] (axis cs:2,3) circle [radius=2pt];
\node[font=\scriptsize, anchor=south west] at (axis cs:2,3) {2};
\filldraw[engblack] (axis cs:4,1.05) circle [radius=2pt];
\node[font=\scriptsize, anchor=south west] at (axis cs:4,1.05) {3};
\filldraw[engblack] (axis cs:2,2.1) circle [radius=2pt];
\node[font=\scriptsize, anchor=north east] at (axis cs:2,2.1) {4};
% labels
\node[engred, font=\scriptsize, anchor=west] at (axis cs:1.45,4.2) {$T_h$};
\node[engblue, font=\scriptsize, anchor=west] at (axis cs:2.7,1.55) {$T_c$};
\end{axis}
\end{tikzpicture}
\caption[The Carnot cycle on a $p$-$V$ diagram]{The Carnot cycle: isothermal
expansion at $T_h$ ($1\to2$, heat $Q_h$ absorbed), adiabatic expansion
($2\to3$, gas cools to $T_c$), isothermal compression at $T_c$ ($3\to4$, heat
$Q_c$ rejected), and adiabatic compression ($4\to1$, gas warms back to $T_h$).
The enclosed area is the net work delivered per cycle. Because the two heat
exchanges happen at the single temperatures $T_h$ and $T_c$, the entropy
transferred in equals the entropy transferred out, $Q_h/T_h = Q_c/T_c$, and the
efficiency reaches the Carnot bound $1 - T_c/T_h$. No cycle operating between
the same two temperatures can do better.}
\label{fig:vol04ch01:carnot-pv}
\end{figure}


% Figure: the Carnot cycle on a T-S diagram, where it is a rectangle.
\begin{figure}[htbp]
\centering
\begin{tikzpicture}
\begin{axis}[
  width=10cm, height=7.5cm,
  xlabel={entropy $S$ (arb.)}, ylabel={temperature $T$ (arb.)},
  xmin=0.5, xmax=4.5, ymin=0, ymax=6,
  axis lines=left,
  tick label style={font=\scriptsize},
  label style={font=\small},
  clip=false,
]
% Carnot rectangle: T_h=5, T_c=2, S from 1 to 4.
% 1->2 isothermal expansion at T_h: S increases 1->4
\addplot[engred, very thick] coordinates {(1,5) (4,5)};
% 2->3 adiabatic (isentropic) expansion: S=4 constant, T 5->2
\addplot[englime!50!black, very thick] coordinates {(4,5) (4,2)};
% 3->4 isothermal compression at T_c: S decreases 4->1
\addplot[engblue, very thick] coordinates {(4,2) (1,2)};
% 4->1 adiabatic compression: S=1 constant, T 2->5
\addplot[engorange, very thick] coordinates {(1,2) (1,5)};
% shade the enclosed area
\addplot[fill=engblue!10, draw=none] coordinates {(1,2) (4,2) (4,5) (1,5)} \closedcycle;
% Qh = area under top edge down to T=0
\node[engred, font=\scriptsize] at (axis cs:2.5,5.3) {$Q_h = T_h \Delta S$};
\node[engblue, font=\scriptsize] at (axis cs:2.5,1.6) {$Q_c = T_c \Delta S$};
\node[font=\scriptsize] at (axis cs:2.5,3.5) {net work $= (T_h-T_c)\Delta S$};
% state labels
\filldraw[engblack] (axis cs:1,5) circle [radius=2pt];
\node[font=\scriptsize, anchor=south east] at (axis cs:1,5) {1};
\filldraw[engblack] (axis cs:4,5) circle [radius=2pt];
\node[font=\scriptsize, anchor=south west] at (axis cs:4,5) {2};
\filldraw[engblack] (axis cs:4,2) circle [radius=2pt];
\node[font=\scriptsize, anchor=north west] at (axis cs:4,2) {3};
\filldraw[engblack] (axis cs:1,2) circle [radius=2pt];
\node[font=\scriptsize, anchor=north east] at (axis cs:1,2) {4};
\end{axis}
\end{tikzpicture}
\caption[The Carnot cycle on a $T$-$S$ diagram]{The same Carnot cycle drawn on
temperature-entropy axes is a rectangle. Heat exchanged along any reversible
path is the area under the path, $Q = \int T\,dS$, so the heat absorbed at the
top edge is $Q_h = T_h\,\Delta S$ and the heat rejected at the bottom edge is
$Q_c = T_c\,\Delta S$. The two adiabatic legs are vertical because a reversible
adiabatic process holds entropy constant. The enclosed rectangle is the net
work $(T_h - T_c)\,\Delta S$, and the efficiency $1 - T_c/T_h$ reads directly
off the ratio of the two horizontal edges' heights.}
\label{fig:vol04ch01:carnot-ts}
\end{figure}


\begin{estimation}
\textbf{How much electrical work does a household refrigerator consume
per day to maintain $4\,\degC$ inside a typical kitchen at
$22\,\degC$?}

\smallskip
\textit{Estimate, before reading on.}

\smallskip
A working estimate: the refrigerator must reject the heat it pumps out
plus the work input to its compressor. The heat to be pumped is the
sum of conductive heat leak through the cabinet walls, infiltration
when the door opens, and the heat of warm items placed inside.

The cabinet has internal volume roughly $0.5\,\text{m}^3$ and surface
area $A \sim 3\,\text{m}^2$. The insulation is polyurethane foam, $5\,
\si{\centi\meter}$ thick, with thermal conductivity $k \approx
0.025\,\text{W}/(\text{m}\cdot\text{K})$. The temperature drop across
the wall is $18\,\si{\kelvin}$. The conduction heat leak is
$\dot{Q}_{\text{leak}} \approx k A \Delta T / L = 0.025 \times 3
\times 18 / 0.05 = 27\,\text{W}$.

Door openings and warm items add intermittently; estimate an average
contribution comparable to the steady leak, so $\dot{Q}_{\text{cold}}
\approx 50\,\text{W}$.

A household refrigerator running on a vapour-compression cycle with
condenser at $\sim 35\,\degC = 308\,\si{\kelvin}$ and evaporator at
$\sim -5\,\degC = 268\,\si{\kelvin}$ has an ideal Carnot coefficient of
performance $\text{COP}_{\text{Carnot}} = T_c/(T_h - T_c) = 268/40
\approx 6.7$. A real refrigerator achieves perhaps $40\,\%$ of the
Carnot bound, so $\text{COP}_{\text{actual}} \approx 2.5$ to $3$. The
required electrical power is $\dot{W} = \dot{Q}_{\text{cold}}/\text{COP}
\approx 50/2.7 \approx 18\,\text{W}$.

Over a $24\,\text{h}$ day, the energy consumption is $18 \times 24
\approx 0.43\,\text{kWh}/\text{day}$, or roughly $160\,\text{kWh}/
\text{year}$.

Manufacturers' specifications for ENERGY STAR refrigerators of this
size cluster between $300$ and $450\,\text{kWh}/\text{year}$
\cite{gen:mackay2009}. Our estimate is a factor of two low; the
discrepancy is dominated by larger units, taller temperature
differences in real cycles, defrost cycles, and warmer ambient.

The postmortem: the estimate was honest about the conduction term and
fragile on the door-opening contribution and the cycle efficiency. The
working calculation gets the order of magnitude (tens of watts of
electrical demand) right, which is what the next-stage design
discussion needs.
\end{estimation}

The refrigerator pumps heat $Q_c$ out of the cold cabinet and the
compressor adds work $W$; the condenser then rejects the sum
$Q_h = Q_c + W$ to the kitchen, as
Figure~\ref{fig:vol04ch01:refrigerator-sankey} draws to scale for the
estimate above. The script \texttt{code/refrigerator\_cop.py} reads a
week of logged consumption and cabinet temperatures from
\texttt{data/refrigerator\_power.csv} and returns the measured
coefficient of performance against the Carnot bound; it is the
analysis engine for the chapter project. Property data for the working
gases of this chapter is collected in \texttt{data/gas\_properties.csv}.

% Figure: energy-flow (Sankey-style) diagram for a refrigerator. Qc drawn from
% cold space + W input combine into Qh rejected at the condenser.
\begin{figure}[htbp]
\centering
\begin{tikzpicture}[font=\small]
  % cold space heat in (Qc ~ 50 W): band height 1.0
  \fill[engblue!35] (0,3.0) rectangle (3.0,4.0);
  \node[font=\scriptsize, anchor=east] at (-0.1,3.5) {$Q_c \approx 50\,\text{W}$};
  \node[engblue, font=\scriptsize] at (1.5,3.5) {from cold space};
  % work input (W ~ 18 W): band height 0.36
  \fill[englime!45] (0,2.4) rectangle (3.0,2.76);
  \node[font=\scriptsize, anchor=east] at (-0.1,2.58) {$W \approx 18\,\text{W}$};
  \node[font=\scriptsize] at (1.5,2.58) {compressor};
  % merge into Qh (height 1.36) leaving to the right at condenser
  \fill[engred!35] (3.0,2.4) -- (5.2,2.4) -- (5.2,4.0) -- (3.0,4.0) -- cycle;
  \node[engred, font=\scriptsize] at (4.1,3.2) {$Q_h \approx 68\,\text{W}$};
  \node[engred, font=\scriptsize, anchor=west] at (5.3,3.2) {rejected at condenser};
  % boundary marker
  \draw[engblack, very thick, dashed] (3.0,1.8) -- (3.0,4.4);
  \node[enggray, font=\scriptsize, anchor=north] at (3.0,1.8) {refrigerator};
\end{tikzpicture}
\caption[Energy flow through a household refrigerator]{Energy flows through a
refrigerator at steady state. The unit pumps heat $Q_c$ out of the cold cabinet
and the compressor adds work $W$; the first law forces the condenser to reject
the sum $Q_h = Q_c + W$ to the kitchen. A refrigerator therefore warms the room
it stands in, by more than the electrical power it draws. The widths are drawn
to scale for the worked estimate of roughly $50\,\text{W}$ pumped, $18\,\text{W}$
of electrical input, and $68\,\text{W}$ rejected; the coefficient of
performance $Q_c/W \approx 2.8$ is the ratio of the two left-hand bands.}
\label{fig:vol04ch01:refrigerator-sankey}
\end{figure}


\section{Failure: perpetual motion machines and how to recognise them in proposals}\pathtag{core}

A perpetual-motion machine is a proposed device that violates one of
the laws of thermodynamics. The discipline of recognising such
proposals is part of every working engineer's tool kit, because
proposals that violate the laws continue to appear in engineering
journals, patent filings, venture pitches, and consumer products. The
proposal is sometimes naive, sometimes deliberate, and sometimes the
output of someone who has confused themselves with a long calculation.
The reader who can name the law violated and the mechanism by which
the proposal violates it is protected against most of these.

\engterm{Perpetual motion of the first kind} violates the first law:
the machine produces net work without an energy input. Any closed
machine that runs forever while delivering useful work is a
first-kind perpetual motion machine. Examples that recur in the
literature:

\begin{itemize}[leftmargin=*]
  \item The over-balanced wheel: a wheel with weights that slide
    outward on one side and inward on the other, claimed to produce a
    net torque. The wheel's centre of mass is fixed by its geometry; a
    careful balance shows the moments cancel.
  \item Magnetic motors that allegedly run on the field of permanent
    magnets without input. The energy stored in a permanent-magnet
    field is finite; once extracted, the magnet is demagnetised.
  \item Water-wheels or weight-driven devices that claim to lift more
    weight than they descend with. The work output cannot exceed the
    gravitational potential energy descended through.
  \item Devices proposed in patent filings as ``over-unity'': the
    measured output exceeds the measured input. The discrepancy is
    always traced to a measurement error, a hidden input (battery,
    grid connection, chemical fuel), or a misidentified flow.
\end{itemize}

The diagnostic is simple: draw the boundary, list every energy flow
across it, and verify the balance. A claim that the input is less
than the output points to an uncounted input, not to a law violation.

\engterm{Perpetual motion of the second kind} violates the second
law: the machine converts heat to work without rejecting heat to a
colder reservoir, or it transfers heat from a colder body to a hotter
body without input work. The class includes more sophisticated
proposals than the first kind, because the second-law violation is
not always obvious from a single balance. Examples:

\begin{itemize}[leftmargin=*]
  \item A machine that extracts work from the heat of a single
    reservoir (the ocean, the atmosphere, the warm earth). The first
    law is satisfied, but the Kelvin-Planck statement is violated: a
    cyclic device cannot extract work from a single reservoir.
  \item Maxwell's demon: a microscopic agent that sorts fast and slow
    molecules to produce a temperature difference, allowing work
    extraction. The resolution (Landauer, Bennett) is that the
    demon's measurement and memory-erasure costs are at least as much
    work as the demon can extract; the second law applies to the
    combined demon-plus-gas system.
  \item Self-cooling devices that allegedly require no input work to
    maintain a temperature difference. The Clausius statement is
    violated; the device, if real, would equilibrate to the ambient
    temperature.
  \item Vapour-compression cycles that claim a coefficient of
    performance exceeding the Carnot bound $T_c/(T_h - T_c)$ for
    their operating temperatures. The Carnot bound is a theorem of
    thermodynamics; a higher COP is a measurement error or a
    misidentification of the operating temperatures.
\end{itemize}

\engterm{Perpetual motion of the third kind} violates the third law:
the machine reaches absolute zero in a finite number of cooling
stages. The class is rarely encountered in engineering proposals; it
is included in some classifications and not in others. The third-law
violation is the operational form of the unreachability of $T = 0$.

The discipline for evaluating a proposed energy device is mechanical:

\begin{enumerate}[leftmargin=*]
  \item Draw the system boundary explicitly, on paper.
  \item List every flow of mass, work, and heat across the boundary,
    with sign convention chosen and stated.
  \item Apply the first law: does $\Delta E = Q + W$ close, with all
    inputs counted? If the claim is over-unity at this stage, the
    missing input is the diagnostic.
  \item Apply the second law: does the proposed cycle satisfy the
    Clausius and Kelvin-Planck statements? Compute the entropy
    generation and confirm it is non-negative.
  \item Compute the Carnot bound for any heat engine or refrigeration
    claim and confirm the proposed efficiency or COP lies below the
    bound.
  \item If the proposal involves chemical or electrochemical
    transformations, verify the Gibbs free-energy balance (the
    maximum useful work is $-\Delta G$ at constant $T$ and $p$).
\end{enumerate}

A proposal that survives the six steps is not a perpetual motion
machine. A proposal that fails any step is, regardless of how
sophisticated the rest of the apparatus appears. The history of
patent offices is filled with rejected over-unity claims; the U.S.
Patent and Trademark Office's standing rule that perpetual-motion
patents require a working prototype before examination is an
operational acknowledgment of the recurrence of the failure mode.

\begin{principle}[The three laws as filter]
A proposed energy device that violates the first law produces work
from no input; a proposal that violates the second law produces work
from a single reservoir or moves heat against the gradient without
input work; a proposal that violates the third law reaches absolute
zero. Every such proposal can be diagnosed by drawing the boundary,
listing the flows, and applying the three balance equations (energy,
entropy, third-law zero) in turn. The diagnostic survives every
generation of perpetual-motion proposal because the laws are
empirical, not contingent on the proposer's intent.
\end{principle}

\section{Project}\pathtag{core}

\begin{project}[Instrument a household refrigerator]
The reader instruments a working household refrigerator over one week
and measures its work input, its heat rejection, and its coefficient
of performance. The deliverable is a fifteen- to twenty-five-page
report.

\textbf{Track}. Household instrumentation plus analysis.
\textbf{Hazard class}: low.

\textbf{Measurements}.

\begin{itemize}[leftmargin=*]
  \item Electrical energy consumption: a plug-in energy meter (Kill A
    Watt or equivalent) records the cumulative kWh delivered to the
    refrigerator over seven days, with the meter zeroed at the start.
  \item Internal cabinet temperature: a logging thermometer (or a
    smartphone thermometer with manual recording at $30\,\text{min}$
    intervals during waking hours) records the temperature at the
    same location for the full week.
  \item Ambient kitchen temperature: a logging thermometer at the
    refrigerator's external location records ambient over the same
    week.
  \item Condenser-coil temperature: an infrared thermometer or
    surface-contact thermometer reads the coil temperature at four
    times during the week.
  \item Cabinet contents: photograph the cabinet contents at the
    start and end of the week, noting any large additions or
    removals.
\end{itemize}

\textbf{Analysis}.

\begin{itemize}[leftmargin=*]
  \item Plot the internal temperature, ambient temperature, and
    condenser temperature over the week. Identify the
    compressor-on and compressor-off cycles in the temperature trace.
  \item Estimate the heat leak through the cabinet walls from the
    measured temperature difference, the cabinet's measured external
    dimensions, and an assumed insulation thickness and thermal
    conductivity (state the assumptions).
  \item Compute the actual coefficient of performance from $\text{COP}
    = \dot{Q}_{\text{cold}}/\dot{W}$, with $\dot{Q}_{\text{cold}}$
    estimated as the steady-state heat leak and $\dot{W}$ measured
    from the kWh meter divided by the week.
  \item Compare the measured COP to the Carnot bound for the measured
    evaporator and condenser temperatures.
  \item Identify the largest source of entropy generation in the
    cycle (the finite-temperature heat exchange at the condenser, the
    throttle valve, the compressor's mechanical losses, or the
    cabinet's heat leak), with a one-paragraph defence.
\end{itemize}

\textbf{Reflection ($1{,}500$ to $2{,}500$ words)}. The reader
addresses:

\begin{itemize}[leftmargin=*]
  \item What fraction of the Carnot bound did the refrigerator
    achieve? What does the gap come from?
  \item Which design changes would close the gap, and at what
    additional manufacturing or operational cost?
  \item One feature of the refrigerator's operation that the data
    revealed and the reader had not previously known.
\end{itemize}

\textbf{Deliverable}. The data log (CSV or spreadsheet), the
calculations, the report, and a one-paragraph summary suitable for
sharing with a non-engineering household member.

\textbf{Time}. One week of passive monitoring, four to eight hours of
data reduction and analysis, four to six hours of write-up.
\end{project}

\section{Exercises}

\subsection*{Calculation}\pathtag{core}

\begin{exercise}
A closed, rigid vessel contains $2\,\text{kg}$ of air at $1\,\text{atm}$
and $20\,\degC$. The vessel is heated until the pressure reaches $3\,
\text{atm}$. Treating air as an ideal gas with $c_v = 0.718\,
\text{kJ}/(\text{kg}\cdot\si{\kelvin})$, find the heat added, the
final temperature, and the work done.
\paragraph{Solution sketch.} Rigid vessel: $V$ constant, so $W = 0$. The
ideal-gas law at fixed mass and volume gives $T_2/T_1 = p_2/p_1 = 3$, so
$T_2 = 3 \times 293\,\si{\kelvin} = 879\,\si{\kelvin}$
(i.e.\ $606\,\degC$). With $W = 0$ the first law is $Q = \Delta U =
m c_v (T_2 - T_1) = 2 \times 0.718 \times (879 - 293) =
841\,\text{kJ}$. All the added heat becomes internal energy because the
gas does no boundary work.
\end{exercise}

\begin{exercise}
A piston-cylinder device contains $0.5\,\text{kg}$ of nitrogen at
$200\,\text{kPa}$ and $300\,\si{\kelvin}$. The gas is compressed
isothermally to $800\,\text{kPa}$. Find the work done on the gas, the
heat removed, and the entropy change of the gas.
\paragraph{Solution sketch.} Nitrogen has $R = 8.314/0.028 =
296.8\,\text{J}/(\text{kg}\cdot\si{\kelvin})$. Isothermal, so
$\Delta U = 0$ and $Q = -W$. The volume ratio equals the inverse
pressure ratio, $V_2/V_1 = p_1/p_2 = 1/4$. Work on the gas:
$W = -mRT \ln(V_2/V_1) = -0.5 \times 296.8 \times 300 \times \ln(1/4)
= +61.7\,\text{kJ}$. Heat removed: $Q = -W = -61.7\,\text{kJ}$ (rejected
to the reservoir). Entropy change of the gas:
$\Delta S = mR \ln(V_2/V_1) = 0.5 \times 296.8 \times \ln(0.25) =
-205.7\,\text{J}/\si{\kelvin}$. The gas entropy falls; the reservoir's
rises by the same amount, so the total is zero for the reversible
idealisation.
\end{exercise}

\begin{exercise}
One kilogram of air ($\gamma = 1.4$, $R = 287\,\text{J}/(\text{kg}
\cdot\si{\kelvin})$) is compressed adiabatically and reversibly from
$100\,\text{kPa}$ and $300\,\si{\kelvin}$ to $500\,\text{kPa}$. Find
the final temperature, the work done on the gas, and the entropy
change.
\paragraph{Solution sketch.} Reversible adiabatic, so use the
pressure-temperature relation
$T_2 = T_1 (p_2/p_1)^{(\gamma-1)/\gamma} = 300 \times 5^{0.4/1.4} =
300 \times 5^{0.2857} = 475\,\si{\kelvin}$. With
$c_v = R/(\gamma-1) = 287/0.4 = 717.5\,\text{J}/(\text{kg}\cdot
\si{\kelvin})$, the work on the gas equals the internal-energy rise
($Q = 0$): $W = m c_v (T_2 - T_1) = 1 \times 717.5 \times (475 - 300)
= 126\,\text{kJ}$. Reversible and adiabatic means $\Delta S = 0$.
\end{exercise}

\begin{exercise}
Steam enters a turbine at $3\,\text{MPa}$ and $400\,\degC$ with
enthalpy $h_1 = 3231\,\text{kJ}/\text{kg}$ and exits at $50\,\text{kPa}$
with enthalpy $h_2 = 2403\,\text{kJ}/\text{kg}$. The mass flow rate is
$5\,\text{kg}/\text{s}$. Find the shaft power output, assuming
negligible heat loss and negligible kinetic- and potential-energy
changes.
\paragraph{Solution sketch.} Steady-state open system, adiabatic, KE
and PE negligible: the balance reduces to
$\dot{W}_s = \dot{m}(h_2 - h_1)$ for work on the CV, so the power
delivered by the turbine is $\dot{m}(h_1 - h_2) = 5 \times (3231 -
2403) = 5 \times 828 = 4140\,\text{kW} = 4.14\,\text{MW}$. The shaft
work on the CV is $-4.14\,\text{MW}$; the turbine extracts
$4.14\,\text{MW}$. The check is reproduced by
\texttt{code/first\_law\_balance.py}.
\end{exercise}

\begin{exercise}
Water enters an electric water heater at $15\,\degC$ at a flow rate of
$10\,\text{L}/\text{min}$ and leaves at $60\,\degC$. Find the
electrical power required, assuming all input power becomes heat in
the water and the water's specific heat is $4.18\,\text{kJ}/(\text{kg}
\cdot\si{\kelvin})$.
\paragraph{Solution sketch.} Mass flow:
$\dot{m} = 10\,\text{L}/\text{min} \times 1\,\text{kg}/\text{L} =
0.167\,\text{kg}/\text{s}$. Required heat rate:
$\dot{Q} = \dot{m} c \Delta T = 0.167 \times 4.18 \times (60 - 15) =
31.4\,\text{kW}$. With all electrical input converted to heat in the
water, the electrical power equals $\dot{Q} \approx 31\,\text{kW}$.
(A real domestic supply rarely delivers this on a single circuit,
which is why instantaneous heaters of this duty are uncommon.)
\end{exercise}

\begin{exercise}
A heat pump operates between an outdoor reservoir at $0\,\degC$ and an
indoor reservoir at $22\,\degC$. The Carnot coefficient of performance
for heating is $\text{COP}_h = T_h/(T_h - T_c)$. Compute the Carnot
COP. If the actual heat pump achieves $50\,\%$ of the Carnot value,
how many watts of electrical input are needed to deliver $5\,\text{kW}$
of heating to the indoor space?
\paragraph{Solution sketch.} Convert: $T_h = 295\,\si{\kelvin}$,
$T_c = 273\,\si{\kelvin}$. Carnot heating COP:
$\text{COP}_h = T_h/(T_h - T_c) = 295/22 = 13.4$. Actual COP at half
the bound: $6.7$. Electrical input to deliver $5\,\text{kW}$ of heat:
$\dot{W} = \dot{Q}_h/\text{COP} = 5000/6.7 = 746\,\text{W}$. A
resistance heater would draw the full $5\,\text{kW}$; the heat pump
delivers the same heat for roughly one-seventh of the electricity
because it moves heat rather than generating it.
\end{exercise}

\begin{exercise}
Two streams of water mix in an open vessel: $1\,\text{kg}/\text{s}$ at
$80\,\degC$ and $2\,\text{kg}/\text{s}$ at $20\,\degC$. Find the
outlet temperature and the entropy generation rate in the mixing
process. Assume negligible heat loss to the surroundings.
\paragraph{Solution sketch.} Energy balance on the adiabatic mixer:
$\dot{m}_1 c T_1 + \dot{m}_2 c T_2 = (\dot{m}_1 + \dot{m}_2) c T_{out}$,
so the outlet is the mass-weighted mean
$T_{out} = (1 \times 80 + 2 \times 20)/3 = 40\,\degC = 313\,\si{\kelvin}$.
Entropy generation (no heat crosses the boundary, so
$\dot{S}_{\text{gen}}$ is the net rise in stream entropy), with
$c = 4.18\,\text{kJ}/(\text{kg}\cdot\si{\kelvin})$ and absolute
temperatures $T_1 = 353$, $T_2 = 293$, $T_{out} = 313\,\si{\kelvin}$:
\[
\dot{S}_{\text{gen}} = \dot{m}_1 c \ln\frac{T_{out}}{T_1} +
\dot{m}_2 c \ln\frac{T_{out}}{T_2}
= 4.18\!\left(1\ln\tfrac{313}{353} + 2\ln\tfrac{313}{293}\right).
\]
The first term is $-0.503$, the second $+0.552$, summing to
$+0.049\,\text{kW}/\si{\kelvin} = 49\,\text{W}/\si{\kelvin}$. The
generation is positive, as the second law requires for an irreversible
mixing.
\end{exercise}

\begin{exercise}
A rigid tank of volume $0.1\,\text{m}^3$ initially contains air at
$200\,\text{kPa}$ and $25\,\degC$. A valve is opened and air leaks to
the atmosphere ($100\,\text{kPa}$, $25\,\degC$) until the tank reaches
atmospheric pressure. Find the entropy change of the air that
remained in the tank (assume the leak is slow and isothermal).
\paragraph{Solution sketch.} Track the gas that stays behind. Slowly
and isothermally, it expands to fill the whole tank as the rest leaks
out; its own pressure falls from $200$ to $100\,\text{kPa}$, which for a
control mass at constant $T$ is an isothermal expansion with
$\Delta S = m_{\text{rem}} R \ln(p_1/p_2)$. The final mass in the tank
is $m_{\text{rem}} = p_2 V/(RT) = 100{,}000 \times 0.1/(287 \times
298) = 0.117\,\text{kg}$, so $\Delta S = 0.117 \times 287 \times
\ln 2 = +23.2\,\text{J}/\si{\kelvin}$. The entropy of the retained gas
rises because it expands; the heat to keep it isothermal came from the
surroundings.
\end{exercise}

\begin{exercise}
Compute the entropy change when $1\,\text{mol}$ of nitrogen at
$1\,\text{atm}$ and $25\,\degC$ is mixed with $1\,\text{mol}$ of oxygen
at the same conditions, with the final mixture occupying twice the
initial volume of each. State whether the process is reversible.
\paragraph{Solution sketch.} Each gas doubles the volume available to
it, so the mixing entropy is the sum of two isothermal expansions:
$\Delta S = -R(n_{N_2}\ln x_{N_2} + n_{O_2}\ln x_{O_2})$ with
$x_{N_2} = x_{O_2} = 0.5$. Then $\Delta S = -R(1\ln 0.5 + 1\ln 0.5) =
2R\ln 2 = 2 \times 8.314 \times 0.693 = 11.5\,\text{J}/\si{\kelvin}$.
The process is irreversible: $\Delta S > 0$ for an isolated system, and
unmixing would require external work.
\end{exercise}

\begin{exercise}
A solid copper block ($m = 2\,\text{kg}$, $c = 385\,\text{J}/(\text{kg}
\cdot\si{\kelvin})$) at $100\,\degC$ is dropped into a large lake at
$15\,\degC$. Find the entropy change of the copper, the entropy
change of the lake, and the total entropy change.
\paragraph{Solution sketch.} Absolute temperatures: copper $373\to
288\,\si{\kelvin}$, lake fixed at $288\,\si{\kelvin}$. Copper:
$\Delta S_{Cu} = m c \ln(T_2/T_1) = 2 \times 385 \times \ln(288/373) =
-199\,\text{J}/\si{\kelvin}$. Heat dumped into the lake:
$Q = m c (T_1 - T_2) = 2 \times 385 \times 85 = 65.5\,\text{kJ}$,
received reversibly at the lake's fixed temperature, so
$\Delta S_{\text{lake}} = +Q/T_{\text{lake}} = 65{,}450/288 =
+227\,\text{J}/\si{\kelvin}$. Total:
$\Delta S_{\text{tot}} = -199 + 227 = +28\,\text{J}/\si{\kelvin} > 0$.
The drop is irreversible: heat crossed a finite temperature difference.
\end{exercise}

\subsection*{Derivation}\pathtag{standard}

\begin{exercise}
Starting from the first law for a closed system and the definition of
enthalpy $H = U + pV$, derive $\Delta H = Q$ for a closed system
undergoing a constant-pressure quasi-static process.
\paragraph{Solution sketch.} First law: $\Delta U = Q + W$ with boundary
work $W = -\int p\,dV = -p\,\Delta V$ at constant $p$. So $Q = \Delta U
+ p\,\Delta V$. At constant pressure, $\Delta(pV) = p\,\Delta V$, hence
$Q = \Delta U + \Delta(pV) = \Delta(U + pV) = \Delta H$. The enthalpy
change equals the heat exchanged precisely because the $p\,\Delta V$
flow-work term is absorbed into the definition of $H$.
\end{exercise}

\begin{exercise}
Derive the relation $c_p - c_v = R$ for an ideal gas from the
definitions of $c_p$ and $c_v$ and the ideal-gas equation of state
$pV = nRT$.
\paragraph{Solution sketch.} Per mole, $H = U + pV = U + RT$ for an
ideal gas. Differentiate with respect to $T$:
$dH/dT = dU/dT + R$. The constant-pressure and constant-volume specific
heats are $c_p = (dH/dT)_p$ and $c_v = (dU/dT)_v$; for an ideal gas $U$
depends on $T$ alone, so both derivatives are ordinary, giving
$c_p = c_v + R$, i.e.\ $c_p - c_v = R$. The relation says the extra heat
a constant-pressure process needs over a constant-volume one, per
degree, is exactly the work $R\,dT$ done in expanding.
\end{exercise}

\begin{exercise}
For an ideal gas with constant $c_v$, derive the expression $\Delta S
= n c_v \ln(T_2/T_1) + nR \ln(V_2/V_1)$ from the combined first and
second laws $dU = T\,dS - p\,dV$.
\paragraph{Solution sketch.} Solve the combined law for $dS$:
$dS = (dU + p\,dV)/T$. For an ideal gas $dU = n c_v\,dT$ and $p/T =
nR/V$, so $dS = n c_v\,dT/T + nR\,dV/V$. Integrate each term from state
1 to state 2 with $c_v$ constant:
$\Delta S = n c_v \ln(T_2/T_1) + nR \ln(V_2/V_1)$. The temperature
term is the entropy of heating; the volume term is the entropy of
expansion. The $(T,p)$ form follows by substituting $V \propto T/p$.
\end{exercise}

\begin{exercise}
Derive the Maxwell relation $(\partial S/\partial V)_T = (\partial p/
\partial T)_V$ from the differential of the Helmholtz free energy
$dA = -S\,dT - p\,dV$.
\paragraph{Solution sketch.} From $dA = -S\,dT - p\,dV$ read off the
first partials: $(\partial A/\partial T)_V = -S$ and $(\partial A/
\partial V)_T = -p$. Because $A$ is a state function, its mixed second
partials are equal: $\partial^2 A/\partial V\,\partial T = \partial^2 A/
\partial T\,\partial V$. Differentiating $-S$ with respect to $V$ and
$-p$ with respect to $T$ and equating gives $-(\partial S/\partial V)_T
= -(\partial p/\partial T)_V$, hence $(\partial S/\partial V)_T =
(\partial p/\partial T)_V$. The Maxwell relation trades an entropy
derivative for one obtainable from $pVT$ data.
\end{exercise}

\begin{exercise}
Show that for an ideal gas undergoing a reversible adiabatic process,
$pV^\gamma$ is constant, where $\gamma = c_p/c_v$.
\paragraph{Solution sketch.} Adiabatic and reversible: $\delta Q = 0$,
so the first law per mole gives $c_v\,dT = -p\,dV$. Differentiate the
ideal-gas law $pV = RT$: $p\,dV + V\,dp = R\,dT$. Eliminate $dT$ using
$dT = -p\,dV/c_v$ and use $R = c_p - c_v$:
$p\,dV + V\,dp = -(c_p - c_v)\,p\,dV/c_v$. Collect terms:
$V\,dp = -(c_p/c_v)\,p\,dV$, i.e.\ $dp/p = -\gamma\,dV/V$. Integrating,
$\ln p + \gamma \ln V = \text{const}$, so $pV^\gamma$ is constant.
\end{exercise}

\begin{exercise}
Derive the steady-state open-system first law from the closed-system
first law applied to a control mass passing through a control volume.
Identify the flow-work term $pv$ that combines with $u$ to form $h$.
\paragraph{Solution sketch.} Take a fixed mass that enters the CV
through the inlet and leaves through the outlet. Pushing a slug of mass
$\delta m$ in at pressure $p_1$ requires the upstream fluid to do work
$p_1 v_1\,\delta m$ on it; pushing the slug out at $p_2$ has the CV do
work $p_2 v_2\,\delta m$ on the downstream fluid. The closed-system
balance for the control mass is $\Delta U = Q + W_{\text{shaft}} +
(p_1 v_1 - p_2 v_2)\delta m$. Grouping the flow-work terms with the
internal energy, $u + pv = h$, gives, at steady state with rate forms,
$0 = \dot{Q} + \dot{W}_s + \dot{m}(h_1 - h_2)$, plus the kinetic and
potential terms if retained. The flow work $pv$ is what promotes $u$ to
$h$ for any stream crossing a boundary at non-zero pressure.
\end{exercise}

\begin{exercise}
Show that the Clausius statement and the Kelvin-Planck statement of
the second law are equivalent: a device that violates either could be
combined with a Carnot engine to produce a device that violates the
other.
\paragraph{Solution sketch.} Show the contrapositive both ways.
Suppose a Kelvin-Planck violator exists: a device that takes $Q$ from
the hot reservoir and produces work $W = Q$ with no rejection. Feed
that work to a Carnot refrigerator pumping $Q_c$ from the cold
reservoir and delivering $Q_c + W$ to the hot one. The combined device
moves net heat $Q_c$ from cold to hot with no external work, a Clausius
violation. Conversely, a Clausius violator (heat $Q_c$ from cold to hot
unaided) can be paired with a Carnot engine that takes $Q_c + W$ from
the hot reservoir, returns exactly $Q_c$ to the cold one, and nets work
$W$; the cold reservoir is unchanged, so the pair extracts work from a
single (hot) reservoir, a Kelvin-Planck violation. Each statement
implies the other.
\end{exercise}

\subsection*{Estimation}\pathtag{core}

\begin{exercise}
Estimate the heat capacity of one kilogram of dry air at constant
pressure at $300\,\si{\kelvin}$, and the heat required to warm a
$30\,\text{m}^3$ room from $15\,\degC$ to $22\,\degC$ without
infiltration.
\paragraph{Reference estimate.} Air $c_p \approx
1.0\,\text{kJ}/(\text{kg}\cdot\si{\kelvin})$ near room temperature
(from $c_p = \gamma R/((\gamma-1)M) = 1.4 \times 8.314/(0.4 \times
0.029) \approx 1.0$). Room air mass: density $\approx
1.2\,\text{kg}/\text{m}^3$ gives $m = 1.2 \times 30 = 36\,\text{kg}$.
Heat to warm by $7\,\si{\kelvin}$: $Q = m c_p \Delta T = 36 \times 1.0
\times 7 = 252\,\text{kJ}$, about $0.07\,\text{kWh}$. A few minutes of a
$1\,\text{kW}$ heater, ignoring the much larger task of warming the
walls and contents.
\end{exercise}

\begin{exercise}
Estimate the work required to compress one cubic metre of air from
$1\,\text{atm}$ to $7\,\text{atm}$ isothermally at $300\,\si{\kelvin}$.
Compare to the work for the same compression done adiabatically.
\paragraph{Reference estimate.} Moles in $1\,\text{m}^3$ at $1\,\text{atm}$,
$300\,\si{\kelvin}$: $n = pV/RT = 101{,}325 \times 1/(8.314 \times 300)
= 40.6\,\text{mol}$. Isothermal work on the gas:
$W = nRT\ln(p_2/p_1) = 40.6 \times 8.314 \times 300 \times \ln 7 =
197\,\text{kJ}$. Adiabatic: $T_2 = 300 \times 7^{0.4/1.4} = 523\,
\si{\kelvin}$, and $W = n c_v \Delta T = 40.6 \times (2.5 \times 8.314)
\times 223 = 188\,\text{kJ}$ of internal-energy rise. The adiabatic
work is comparable but leaves the gas hot; if the gas is then cooled
back to $300\,\si{\kelvin}$ at constant pressure the total energy bill
exceeds the isothermal case. Isothermal compression is the cheaper
route, which is why large compressors are intercooled.
\end{exercise}

\begin{exercise}
Estimate the entropy change of the universe when a kettle boils
$1\,\text{L}$ of water from $20\,\degC$ to $100\,\degC$ on an
electric burner that operates at $700\,\degC$.
\paragraph{Reference estimate.} Heat into the water:
$Q = m c \Delta T = 1 \times 4.18 \times 80 = 334\,\text{kJ}$. Water
entropy gain (heating, not a phase change):
$\Delta S_{\text{water}} = mc\ln(T_2/T_1) = 4180\ln(373/293) =
+1010\,\text{J}/\si{\kelvin}$. The burner surface at $973\,\si{\kelvin}$
gives up the same $334\,\text{kJ}$, losing
$\Delta S_{\text{burner}} = -334{,}000/973 = -343\,\text{J}/
\si{\kelvin}$. Universe:
$\Delta S \approx 1010 - 343 = +670\,\text{J}/\si{\kelvin}$, positive
and dominated by delivering heat across the large burner-to-water
temperature gap.
\end{exercise}

\begin{exercise}
Estimate the order of magnitude of the daily food-energy intake of an
adult human in joules, and compare to the daily energy consumption of
a household refrigerator.
\paragraph{Reference estimate.} A daily intake of $2000\,\text{kcal}$
is $2000 \times 4184 = 8.4\,\text{MJ}$, equivalent to $2.3\,\text{kWh}$.
A refrigerator drawing $\sim 1\,\text{kWh}/\text{day}$ (label figures
of $\sim 350\,\text{kWh}/\text{year}$) uses roughly half the energy a
person eats. Averaged over the day the human metabolises at about
$8.4\,\text{MJ}/86{,}400\,\text{s} \approx 100\,\text{W}$, the same
order as the refrigerator's average electrical draw.
\end{exercise}

\begin{exercise}
Estimate the maximum work that can be obtained from $1\,\text{kg}$ of
hot tea at $80\,\degC$ delivered to a $20\,\degC$ environment, if the
tea is brought to ambient temperature reversibly. State the
calculation as an entropy bound.
\paragraph{Reference estimate.} The maximum work is the exergy of the
hot tea relative to the $T_0 = 293\,\si{\kelvin}$ environment,
$W_{\max} = mc[(T_1 - T_0) - T_0\ln(T_1/T_0)]$ with $T_1 = 353\,
\si{\kelvin}$. Numerically
$W_{\max} = 4180[(353-293) - 293\ln(353/293)] = 4180[60 - 52.6] =
22.6\,\text{kJ}$. The bracketed form is the entropy bound: the
$60\,\si{\kelvin}$ of enthalpy drop is discounted by $T_0$ times the
entropy the tea sheds, $T_0\,\Delta S_{\text{tea}}$, which cannot be
recovered as work. The heat content released is $mc\Delta T =
251\,\text{kJ}$, so only about $9\,\%$ of it is available as work
against this environment.
\end{exercise}

\subsection*{Simulation}\pathtag{mastery}

\begin{exercise}
Implement a $p$-$V$ diagram plotter for an ideal-gas cycle composed of
two isothermal and two adiabatic segments (a Carnot cycle). Plot the
cycle and compute the net work and the thermal efficiency. Verify
that the computed efficiency matches the Carnot bound $1 - T_c/T_h$.
\paragraph{Solution sketch.} Fix $T_h$, $T_c$, $\gamma$, and the start
state. Generate each leg parametrically: isotherms as $p = nRT/V$ over
the chosen $V$ range, adiabats as $p = p_{\text{start}}(V_{\text{start}}
/V)^\gamma$, matching the adiabatic relation $T V^{\gamma-1} =
\text{const}$ to set the corner volumes. Net work is the enclosed area
$\oint p\,dV$, evaluated by the trapezoid rule on the plotted points,
or equivalently from $Q_h - Q_c$ with $Q_h = nRT_h\ln(V_2/V_1)$.
Efficiency $\eta = W_{\text{net}}/Q_h$ should equal $1 - T_c/T_h$ to
plotting precision. The reference implementation is
\texttt{code/carnot\_cycle.py}, which asserts the match to machine
precision.
\end{exercise}

\begin{exercise}
Simulate the mixing of two ideal gases in a $1\,\text{L}$ container
divided by a partition, using a stochastic particle simulation
($10^4$ particles, periodic boundary conditions, elastic collisions).
Plot the time-evolution of the species concentrations on each side and
the total entropy (computed from the bin populations) as a function
of time.
\paragraph{Solution sketch.} Initialise species $A$ on the left and
$B$ on the right; advance positions with elastic collisions (or, more
simply, a diffusive random hop between bins). At each step estimate the
entropy from the bin compositions using $-R\sum_i x_i\ln x_i$ summed
over bins, weighted by occupancy. The concentrations equalise across
the partition and the entropy rises monotonically (apart from
stochastic noise) toward $R\ln 2$ per mole for an equimolar start. The
curve is the kinetic face of the mixing-entropy formula of the worked
examples; \texttt{code/entropy\_of\_mixing.py} carries a minimal
implementation and prints the limiting value.
\end{exercise}

\begin{exercise}
Implement a steady-state open-system energy-balance solver: given
inlet conditions, outlet conditions, work, and heat for any chosen
control volume, return the residual of the first law and flag any
inconsistency. Verify on the steam turbine of the Calculation
exercises.
\paragraph{Solution sketch.} Implement the residual
$r = \dot{Q} + \dot{W}_s + \dot{m}(h_1 - h_2 + \tfrac12(V_1^2 - V_2^2)
+ g(z_1 - z_2))$ and flag $|r|$ above a tolerance. For the steam
turbine ($\dot{m} = 5$, $h_1 = 3231$, $h_2 = 2403\,\text{kJ}/\text{kg}$,
adiabatic, KE and PE negligible), solving for $\dot{W}_s$ returns
$-4.14\,\text{MW}$ and the residual is zero. The reference
implementation \texttt{code/first\_law\_balance.py} carries both the
residual and the solve, with an assertion on the turbine case.
\end{exercise}

\subsection*{Design}\pathtag{standard}

\begin{exercise}
Specify the size of an electric water heater (in kW) to deliver
$50\,\text{L}/\text{min}$ at $45\,\degC$ to a residential shower,
with inlet water at $10\,\degC$. State the assumptions and the
energy efficiency assumed.
\paragraph{Solution sketch.} Flow $\dot{m} = 50/60 =
0.83\,\text{kg}/\text{s}$, rise $\Delta T = 35\,\si{\kelvin}$. Heat
rate $\dot{Q} = \dot{m}c\Delta T = 0.83 \times 4.18 \times 35 =
122\,\text{kW}$. Resistance heating converts nearly all the electrical
input to heat at the element, so specify roughly $120$ to
$130\,\text{kW}$. (This is
implausibly large for a domestic supply, which is the engineering
point: a tankless heater at this flow and rise needs a commercial feed,
so real showers use storage tanks or lower flow.) Assumptions: steady
flow, no distribution loss, water $c = 4.18\,\text{kJ}/(\text{kg}\cdot
\si{\kelvin})$.
\end{exercise}

\begin{exercise}
Design a passive solar water heating panel that delivers
$200\,\text{L}/\text{day}$ at $50\,\degC$ in a climate where the
inlet water is at $15\,\degC$ and the average solar irradiance is
$200\,\text{W}/\text{m}^2$ over the day. Compute the required
collector area and state the heat-loss assumption.
\paragraph{Solution sketch.} Daily heat demand:
$Q = 200 \times 4.18 \times 35 = 29.3\,\text{MJ}/\text{day}$. Daily
insolation per square metre: $200\,\text{W}/\text{m}^2 \times
86{,}400\,\text{s} = 17.3\,\text{MJ}/(\text{m}^2\cdot\text{day})$. With
a perfect collector the area is $29.3/17.3 = 1.7\,\text{m}^2$; with a
realistic collector efficiency of $50\,\%$ (optical and convective
losses), $A \approx 3.4\,\text{m}^2$. State the heat-loss assumption
explicitly: the $50\,\%$ lumps glazing reflection, absorber
re-radiation, and convective loss to ambient; a real datasheet gives
this as the collector's efficiency curve against $(T_{\text{plate}} -
T_{\text{amb}})/G$.
\end{exercise}

\begin{exercise}
Specify a heat-recovery ventilation unit for a residential building
where the indoor temperature is $22\,\degC$, the outdoor temperature
is $-5\,\degC$, and the ventilation rate is $100\,\text{m}^3/\text{h}$.
State the heat-recovery effectiveness target and compute the
auxiliary heating power required after recovery.
\paragraph{Solution sketch.} Mass flow:
$\dot{m} = 1.2 \times 100/3600 = 0.033\,\text{kg}/\text{s}$. Without
recovery, warming the incoming air across $\Delta T = 27\,\si{\kelvin}$
needs $\dot{m}c_p\Delta T = 0.033 \times 1005 \times 27 =
905\,\text{W}$. Set an effectiveness target $\varepsilon = 0.8$
(typical for a good plate or rotary core); the supply air leaves the
exchanger at $-5 + 0.8 \times 27 = 16.6\,\degC$. The auxiliary heater
then supplies $\dot{m}c_p(22 - 16.6) = 0.033 \times 1005 \times 5.4 =
181\,\text{W}$, a fivefold reduction. The recovery does not violate the
second law: it transfers heat from the warm exhaust to the cold supply
across a finite gradient, generating entropy, and never lifts the
supply above the exhaust temperature.
\end{exercise}

\begin{exercise}
A household kettle holds $1.7\,\text{L}$ of water and boils a full
load from $20\,\degC$ to $100\,\degC$. Specify the electrical power
rating to bring the water to a boil in not more than three minutes,
and state the efficiency assumption.
\paragraph{Solution sketch.} Heat to boil:
$Q = mc\Delta T = 1.7 \times 4.18 \times 80 = 569\,\text{kJ}$. Time
limit $t = 180\,\text{s}$, so the heat rate is $569/180 =
3.16\,\text{kW}$ at perfect efficiency. Allow for heat lost to the
kettle body and steam by assuming $85\,\%$ efficiency, giving a rating
of $3.16/0.85 = 3.7\,\text{kW}$. A standard domestic kettle is rated
$2$ to $3\,\text{kW}$, so a full $1.7\,\text{L}$ load on a
$2.2\,\text{kW}$ element takes closer to five minutes; meeting the
three-minute target needs a $3\,\text{kW}$-plus circuit.
\end{exercise}

\subsection*{Diagnosis and reverse engineering}\pathtag{core}

\begin{exercise}
A household refrigerator's measured energy consumption (from the
appliance label) is $400\,\text{kWh}/\text{year}$. The cabinet
volume is $400\,\text{L}$, the average ambient temperature is
$22\,\degC$, and the internal compartment is $4\,\degC$. Compute the
implied average heat leak in watts and infer the cycle's effective
COP. Identify the largest term in the energy balance.
\paragraph{Solution sketch.} Average electrical power:
$400\,\text{kWh}/\text{year} = 400{,}000/8760 = 45.7\,\text{W}$.
Conduction leak for a $400\,\text{L}$ cabinet (surface area
$\sim 4.5\,\text{m}^2$, foam $5\,\si{\centi\meter}$, $k =
0.025\,\text{W}/(\text{m}\cdot\si{\kelvin})$, $\Delta T = 18\,
\si{\kelvin}$): $\dot{Q}_{\text{leak}} = kA\Delta T/L = 0.025 \times
4.5 \times 18/0.05 = 40\,\text{W}$. Treating that as the entire cold
load gives an implied $\text{COP} = 40/45.7 = 0.9$, which is below one
and therefore not physical for a working cycle. The diagnostic
conclusion: the steady conduction leak is not the largest term. Door
openings, defrost heaters, and the fan and control electronics add to
the load and raise the true cold-side heat above the conduction
estimate, so the real cycle COP exceeds one while the label energy
still resolves to $45.7\,\text{W}$ average. The largest term in a
typical cabinet balance is the combined infiltration-plus-defrost load,
not the wall conduction.
\end{exercise}

\begin{exercise}
A consumer-grade thermos flask holds $1\,\text{L}$ of coffee
initially at $90\,\degC$ in a $20\,\degC$ ambient. After eight hours
the coffee is at $60\,\degC$. Estimate the effective $UA$ value of
the flask (heat-transfer coefficient times area) and the
corresponding overall thermal conductance in $\text{W}/\si{\kelvin}$.
\paragraph{Solution sketch.} Lumped cooling:
$T(t) - T_a = (T_0 - T_a)e^{-UA\,t/(mc)}$. With $T_0 = 90$, $T_a = 20$,
$T(8\,\text{h}) = 60\,\degC$, $m = 1\,\text{kg}$, $c =
4180\,\text{J}/(\text{kg}\cdot\si{\kelvin})$, solve for $UA$:
\[
UA = -\frac{mc}{t}\ln\frac{T(t) - T_a}{T_0 - T_a}
= -\frac{4180}{28{,}800}\ln\frac{40}{70} = 0.081\,\text{W}/\si{\kelvin}.
\]
The overall thermal conductance is about $0.08\,\text{W}/\si{\kelvin}$:
a vacuum flask leaks under a tenth of a watt per kelvin, which is why
it holds heat for hours where a ceramic mug ($UA$ of order
$1\,\text{W}/\si{\kelvin}$) cools in minutes.
\end{exercise}

\begin{exercise}
An espresso machine's heating element is rated $1200\,\text{W}$ and
heats $50\,\text{mL}$ of water from $20\,\degC$ to $90\,\degC$ in
$15\,\text{s}$. Compute the apparent efficiency of the heating
process. Identify where the missing energy is going.
\paragraph{Solution sketch.} Heat into the water:
$Q = mc\Delta T = 0.05 \times 4180 \times 70 = 14{,}630\,\text{J}$.
Electrical energy supplied: $1200 \times 15 = 18{,}000\,\text{J}$.
Apparent efficiency $= 14{,}630/18{,}000 = 0.81$. The missing $19\,\%$
goes into heating the boiler block and group head (their thermal mass
soaks up energy that never reaches this shot), into conduction and
convection loss from the hot metal, and into the warm-up transient
before the element reaches steady output. A machine left on to
pre-heat amortises the block-heating term across many shots, raising
the per-shot efficiency.
\end{exercise}

\subsection*{Failure analysis}\pathtag{core}

\begin{exercise}
A 2024 patent application claims a device that uses a permanent
magnet to drive a small generator that recharges its own input
battery while delivering net power to a load. Identify the law
violated, the boundary on which the violation appears, and the
single measurement that would detect it.
\paragraph{Solution sketch.} This is perpetual motion of the first
kind: it claims net power out while recharging its own source, so a
boundary drawn around the whole sealed device shows work leaving with
no energy entering, in violation of $\Delta E = Q + W$. The permanent
magnet stores a finite field energy that demagnetises once extracted;
it is not a perpetual source. The single detecting measurement: meter
the battery's state of charge (or its terminal energy) over a run with
the external load connected. If the device were real the battery would
not deplete; in every such case it does. Equivalently, integrate net
electrical power across the device boundary over time and confirm it is
not positive.
\end{exercise}

\begin{exercise}
A venture pitch describes an ``ambient energy harvester'' that
extracts useful electrical work from the thermal motion of air at
room temperature without any temperature gradient or other driving
potential. Identify the second-law violation, draw the equivalent
Kelvin-Planck construction, and write a single sentence rejecting the
proposal.
\paragraph{Solution sketch.} The harvester claims to take heat from a
single reservoir (room-temperature air, with no second colder
reservoir) and turn it into work in a cycle. That is exactly the
Kelvin-Planck construction: a cyclic device drawing heat $Q$ from one
reservoir at $T$ and delivering work $W = Q$ with no rejection. The
first law is satisfied; the second law is not. Rejection sentence: ``A
cyclic engine cannot convert heat from a single uniform-temperature
reservoir into work, so without a colder sink or another driving
potential the device produces no net power and is rejected on the
Kelvin-Planck statement.''
\end{exercise}

\begin{exercise}
A consumer product is sold as a refrigerator that operates with no
external power input. The cabinet is internally cooled by an
``advanced thermoelectric'' core. Identify the implied claim about
heat flow, the law it violates, and the empirical test that would
detect the violation.
\paragraph{Solution sketch.} A refrigerator with no power input claims
to move heat from the cold cabinet to the warmer room with no
compensating work, which is the Clausius statement violated outright:
heat does not pass from a colder body to a hotter one of itself. A
thermoelectric (Peltier) core still needs electrical input; called on
to run with none, it cannot maintain a temperature difference. The
empirical test: seal the unit with no electrical or chemical input and
log the internal temperature. By the second law it equilibrates to
ambient; a genuine self-cooling claim would hold the cabinet below room
temperature indefinitely, which never happens. Measuring zero input
power while the cabinet stays cold is the contradiction.
\end{exercise}

\subsection*{Judgment}\pathtag{standard}

\begin{exercise}
A colleague proposes to compute the efficiency of a coal-fired power
plant as $\eta = W_{\text{net}}/Q_{\text{in}}$ but to use the heating
value of coal rather than the enthalpy of the combustion products as
$Q_{\text{in}}$. Write a one-paragraph reply identifying the
inconsistency and naming the correct quantity.
\paragraph{Model-answer sketch.} The numerator and denominator must
refer to the same accounting boundary. ``Heating value'' is a defined
quantity (the enthalpy of combustion to a stated reference state, with
the higher heating value assuming the product water condenses and the
lower heating value assuming it leaves as vapour), so the issue is
consistency rather than the wrong kind of number. The correct
denominator is the heat actually released into the cycle's working
fluid, which for a plant venting hot flue gas and uncondensed water is
the lower heating value times the fuel rate, not the higher. Mixing the
higher heating value into $Q_{\text{in}}$ while the cycle never recovers
the condensation enthalpy understates the efficiency; state which
heating value is used and apply it consistently to both the fuel input
and any recovered heat.
\end{exercise}

\begin{exercise}
A student argues that the second law is "really just" the first law
because both express conservation. Write a one-paragraph reply
distinguishing the two laws by what they conserve and by the
asymmetry of allowed processes.
\paragraph{Model-answer sketch.} The first law conserves energy: the
total energy of an isolated system is constant, and any process keeps
the books balanced. The second law conserves nothing; it asserts an
inequality, that the entropy of an isolated system cannot decrease.
Both directions of a heat transfer between hot and cold bodies conserve
energy, so the first law permits the cold body to warm the hot one. The
second law forbids it, because only one direction increases total
entropy. The first law tells what is allowed by the energy balance; the
second law selects which of the energy-conserving processes actually
happen and sets the maximum work obtainable. They are different
statements about different quantities, not two versions of
conservation.
\end{exercise}

\begin{exercise}
A facility manager argues that the refrigerator should be left fully
loaded because a fuller refrigerator uses less energy. Identify
whether the claim is correct and state the dominant mechanism.
\paragraph{Model-answer sketch.} The claim is partly correct but for
the right reason, and only in a particular operating pattern. At steady
state the compressor work tracks the cabinet heat load, which is set by
the wall conduction and the door-opening losses, not by how full the
cabinet is; a full and an empty refrigerator at the same temperature
leak the same heat through the same walls. The dominant mechanism that
makes a fuller unit appear more efficient is thermal mass: the food and
liquids hold cold, so each door opening warms the air but the bulk
contents recover the temperature with a smaller compressor duty cycle,
and the cabinet rides through outages and defrosts with less swing.
Filling with water bottles to add mass helps a frequently-opened unit;
it does nothing for a sealed one.
\end{exercise}
